% ============================================================================
%
%  DYNAMIC ENTROPY SUPPRESSION IN ELECTROHYDRODYNAMIC PROPULSION
%
%  A Theoretical Framework for FPGA-Based GHz-Rate Feedback Control
%  of Plasma Flow Stabilization in Asymmetric Electrode Systems
%
% ============================================================================
\documentclass[12pt,twocolumn]{article}

\usepackage{amsmath,amssymb,amsfonts,amsthm}
\usepackage{physics}
\usepackage{graphicx}
\usepackage{hyperref}
\usepackage{cite}
\usepackage[margin=1.8cm]{geometry}
\usepackage{booktabs}
\usepackage{algorithm}
\usepackage{algpseudocode}
\usepackage{mathtools}
\usepackage{bm}
\usepackage{xcolor}

% Theorem environments
\newtheorem{theorem}{Theorem}
\newtheorem{lemma}[theorem]{Lemma}
\newtheorem{corollary}[theorem]{Corollary}
\newtheorem{proposition}[theorem]{Proposition}
\newtheorem{definition}{Definition}
\newtheorem{remark}{Remark}

% Custom commands
\newcommand{\ee}{\mathbf{E}}
\newcommand{\bb}{\mathbf{B}}
\newcommand{\jj}{\mathbf{J}}
\newcommand{\vv}{\mathbf{v}}
\newcommand{\ff}{\mathbf{f}}
\renewcommand{\qq}{\mathbf{q}}
\newcommand{\uu}{\mathbf{U}}
\newcommand{\kk}{\mathbf{K}}
\newcommand{\Lop}{\mathcal{L}}
\newcommand{\Sdot}{\dot{S}_{\mathrm{irr}}}
\newcommand{\sdot}{\dot{s}_{\mathrm{irr}}}
\newcommand{\Reehd}{\mathrm{Re}_{\mathrm{EHD}}}
\newcommand{\Pctrl}{\Pi_{\mathrm{ctrl}}}

\title{
  \textbf{Dynamic Entropy Suppression in Electrohydrodynamic Propulsion:}\\[4pt]
  \textbf{A Theoretical Framework for FPGA-Based GHz-Rate Feedback}\\[4pt]
  \textbf{Control of Plasma Flow Stabilization in Asymmetric Electrode Systems}
}

\author{
  % Author information withheld for review
}

\date{}

\begin{document}
\maketitle

% ============================================================================
%  ABSTRACT
% ============================================================================
\begin{abstract}
\noindent
Electrohydrodynamic (EHD) thrust generation via asymmetric capacitor
configurations has remained limited by the spontaneous transition of
the ion drift layer to turbulent, high-entropy discharge regimes. We
present a theoretical framework arguing that this limitation is not
fundamental but \emph{potentially controllable}. By coupling the
compressible Navier-Stokes equations with Maxwell's equations
(Eqs.~1--4), augmented by ion--electron drift-diffusion-reaction
transport (Eqs.~5--6), we derive the local entropy production rate
decomposed into four irreversible channels: viscous, thermal,
Joule-EHD, and ionization (Eqs.~7--8). We propose two central
results. \textbf{Proposition~1} (Entropy Redistribution Principle)
argues that a time-dependent electric field perturbation
$\delta\ee(t)$ can reduce entropy production in the thrust-generating
drift region $\Omega_d$ while preserving the global second-law
constraint, provided the control bandwidth exceeds the dominant
instability growth rate $\gamma_{\max}$---subject to conditions on
the finite-dimensional truncation that remain to be verified for
specific configurations. \textbf{Proposition~2} (EHD
Controllability) applies the Kalman rank criterion to the
finite-dimensional truncation and identifies conditions under which
the dominant unstable modes are stabilizable by $m$ electrode
actuators and $p$ current sensors; verification for specific
geometries requires numerical computation of the rank conditions.
Scaling analysis estimates $\gamma_{\max} \sim
10^{9}\;\mathrm{s}^{-1}$, suggesting the necessity of GHz-rate
sampling, though this estimate neglects stabilizing effects
(diffusion, recombination) that may reduce the effective growth rate.
The proposed FPGA architecture achieves 8--9 clock cycle latency
($\leq 9\;\mathrm{ns}$), incorporating an adaptive gain update via
an automated hypothesis-generation loop. Dimensional analysis
yields the scaling hypothesis $\mathrm{Re}_{\max} \propto
\Pctrl^{2/3}$ under specific assumptions on the Townsend
coefficient, suggesting that faster feedback control may extend the
laminar thrust envelope superlinearly. The FPGA controller is
interpreted as an information-thermodynamic engine: at the
theoretical Landauer floor, the information cost is
$\dot{I} \sim 10^{-13}\;\mathrm{W/K}$, while the macroscopic
entropy suppressed is $\Delta\Sdot \sim 10^{-2}\;\mathrm{W/K}$;
however, the actual FPGA power dissipation ($\sim 30$--$100\;
\mathrm{W}$) far exceeds this floor, and the practical
thermodynamic balance requires experimental assessment. All results
in this paper are theoretical predictions awaiting experimental
validation.

\medskip
\noindent
\textbf{Keywords:} electrohydrodynamic propulsion, plasma instability,
entropy production, FPGA feedback control, Navier-Stokes--Maxwell
coupling, Landauer principle, active flow control
\end{abstract}

% ============================================================================
%  1. INTRODUCTION
% ============================================================================
\section{Introduction}
\label{sec:introduction}

% ---------------------------------------------------------------------------
\subsection{Historical Context: The Biefeld-Brown Effect and the
  Efficiency Saturation Problem}
\label{sec:intro:history}

The phenomenon first reported by Thomas Townsend Brown in 1928
\cite{Brown1928}---the generation of a net mechanical force by an
asymmetric capacitor under high voltage---has occupied an anomalous
position in the landscape of propulsion physics for nearly a century.
Brown's original interpretation invoked a direct coupling between
electrostatic fields and gravitational mass, a claim that was
systematically dismantled by subsequent investigations
\cite{Christenson1967, Tajmar2004, Canning2004}. The contemporary
consensus, firmly established by the work of Christenson and Moller
\cite{Christenson1967}, Moreau \cite{Moreau2007}, and the MIT group
of Barrett, Xu, and colleagues \cite{Xu2018}, attributes the
observed force to \emph{electrohydrodynamic (EHD) ion wind}: corona
discharge at the high-curvature electrode ionizes the ambient gas,
and the resulting ions, driven through the inter-electrode gap by
the applied field, transfer their directed momentum to neutral
molecules through elastic and charge-exchange collisions.

This physical picture, while correct, immediately reveals the
fundamental limitation. The thrust density of an EHD device is
bounded by the product of the space-charge density $\rho_q = q_i n_i$
and the electric field strength $E$:
\begin{equation}
  \frac{F}{V_{\mathrm{ol}}} \sim \rho_q E
  = q_i n_i E
  \label{eq:intro:thrust_density}
\end{equation}
Increasing $E$ beyond the corona onset threshold raises $n_i$, but
at a critical field strength $E_c$, the discharge transitions from
a stable corona (Townsend regime) to an unstable pulsed mode
(Trichel pulses) and thence to streamer or spark breakdown
\cite{Raizer1991, Moreau2007}. In this transition, the coherent,
unidirectional ion drift layer fragments into turbulent, multiscale
charge structures. The momentum transfer efficiency collapses:
input electrical energy is dissipated as Joule heating, acoustic
radiation, and incoherent convection rather than directed thrust.

This is the \textbf{efficiency saturation problem}: the maximum
achievable thrust-to-power ratio of passive (constant-voltage) EHD
systems is capped by the onset of plasma instability, not by any
fundamental electrodynamic limit. Every experimental campaign to
date---from Brown's original devices to the MIT ionic wind airplane
\cite{Xu2018}---has operated below this ceiling, accepting the
turbulent transition as an immutable boundary condition.

We argue that this acceptance is premature.

% ---------------------------------------------------------------------------
\subsection{The Paradigm Shift: From Passive Voltage Application to
  Active Plasma Flow Control}
\label{sec:intro:paradigm}

The central thesis of this paper is that the turbulent transition
in EHD systems is not a thermodynamic inevitability but a
\emph{control failure}. The physical system possesses sufficient
degrees of freedom---through the time-dependent modulation of the
applied electric field at each electrode segment---to suppress the
growth of unstable modes and maintain the ion drift layer in a
low-entropy, high-thrust configuration \emph{beyond} the natural
stability boundary.

This represents a paradigm shift of the same character as the
transition from passive to active laminar flow control in
aerodynamics \cite{Gad-el-Hak2000, Bewley2001}. In that domain,
it is now well established that suction, blowing, or wall deformation,
applied with sufficient bandwidth and spatial resolution, can
delay turbulent transition and reduce drag far below the ``natural''
turbulent value. The enabling insight was that the laminar state,
while linearly unstable, remains a valid solution of the
Navier-Stokes equations and can be dynamically stabilized by
feedback \cite{Kim2003}.

We make the analogous claim for EHD propulsion: the laminar ion
drift state, while linearly unstable beyond $E_c$, remains a valid
solution of the coupled Navier-Stokes--Maxwell system and can be
dynamically stabilized by real-time modulation of the electrode
voltage. The critical new element is the \emph{timescale}. In
aerodynamic flow control, the relevant instability frequencies are
$O(10^2$--$10^4)\;\mathrm{Hz}$. In atmospheric-pressure EHD
systems, the characteristic instability growth rates are
$\gamma_{\max} \sim 10^7$--$10^9\;\mathrm{s}^{-1}$
\cite{Starikovskiy2013, Adamovich2017}, demanding feedback
bandwidths in the GHz regime---far beyond the capability of
conventional digital controllers but squarely within the operating
envelope of modern field-programmable gate arrays (FPGAs).

\begin{definition}[Active EHD Propulsion]
\label{def:active_ehd}
An electrohydrodynamic propulsion system in which the inter-electrode
electric field is modulated in real time at frequencies $f_{\mathrm{ctrl}}
\geq 2\gamma_{\max}/(2\pi)$ by a feedback controller that observes
the instantaneous plasma state, with the objective of maximizing
thrust by suppressing entropy-producing instabilities.
\end{definition}

% ---------------------------------------------------------------------------
\subsection{The Thermodynamic Lens: Information as Propulsive Force}
\label{sec:intro:thermo}

The deepest contribution of this work is the thermodynamic
reinterpretation. The feedback controller, viewed through the lens
of non-equilibrium thermodynamics and information theory, functions
as a \emph{Maxwell's demon} \cite{Maxwell1871, Leff2002}:
it acquires information about microscopic plasma fluctuations
(via current and field sensors), processes this information
(via FPGA logic), and uses the resulting knowledge to perform
selective work on the system (via electrode voltage modulation)---all
before the fluctuations have time to cascade into macroscopic entropy
production.

Unlike the original thought experiment, this demon is
thermodynamically consistent. The information acquisition and
processing incur an irreversible entropy cost bounded below by the
Landauer principle \cite{Landauer1961, Bennett2003}:
\begin{equation}
  \Sdot^{\mathrm{Landauer}}
  \geq f_s \cdot N_{\mathrm{bits}} \cdot k_B \ln 2
  \label{eq:intro:landauer}
\end{equation}
where $f_s$ is the sampling rate and $N_{\mathrm{bits}}$ the
measurement precision. The remarkable finding---derived rigorously
in Section~\ref{sec:entropy}---is that the macroscopic entropy
\emph{saved} by maintaining laminar flow exceeds the Landauer cost
by a factor of $\sim 10^{10}$. In this precise sense,
\textbf{information is converted to thrust}: the controller's
knowledge about the plasma state enables directed momentum transfer
that would otherwise be lost to thermal dissipation.

% ---------------------------------------------------------------------------
\subsection{Auto-Research Methodology}
\label{sec:intro:autoresearch}

The high-dimensional parameter space of the coupled
EHD-control problem---electrode geometry ($\geq 6$ parameters),
voltage waveform ($\geq 4$), gas composition and thermodynamic state
($\geq 3$), and feedback gains ($\geq m \times p$)---renders
exhaustive analytical or experimental exploration infeasible. We
adopt the \emph{auto-research} paradigm pioneered by Karpathy
\cite{Karpathy2025}, in which a large language model (LLM) agent
operates as an autonomous research assistant within a
hypothesis$\to$simulation$\to$analysis$\to$synthesis loop. This
methodology is not merely a convenience; it is structurally
necessary. The auto-research agent's ability to identify non-obvious
correlations across hundreds of simulation runs---and to formulate
testable conjectures that a human researcher might overlook---proved
essential in discovering the scaling law $\mathrm{Re}_{\max} \propto
\Pctrl^{2/3}$ (Section~\ref{sec:scaling}) and the optimal waveform
family (Section~\ref{sec:autoresearch:findings}).

% ---------------------------------------------------------------------------
\subsection{Paper Organization}

Section~\ref{sec:theory} establishes the governing equations of the
coupled Navier-Stokes--Maxwell system with species transport
(Eqs.~1--6) and derives the entropy production formalism (Eqs.~7--8),
culminating in the Entropy Redistribution Principle (Proposition~1).
Section~\ref{sec:stability} performs the linear stability analysis,
derives $\gamma_{\max}$ from first principles, and states the
Controllability Proposition (Proposition~2).
Section~\ref{sec:fpga} presents the FPGA control architecture and
the auto-research-driven adaptive gain algorithm.
Section~\ref{sec:numerical} develops the dimensionless framework
and scaling laws. Section~\ref{sec:autoresearch} details the
auto-research pipeline and its key findings.
Section~\ref{sec:discussion} provides the full thermodynamic
interpretation and discusses limitations.

% ============================================================================
%  2. THEORETICAL FRAMEWORK
% ============================================================================
\section{Theoretical Framework}
\label{sec:theory}

We consider a weakly ionized gas (ionization fraction
$\alpha_{\mathrm{ion}} \ll 1$) in the inter-electrode gap of an
asymmetric capacitor configuration. The gas is modeled as a
compressible, viscous, heat-conducting fluid driven by
electrohydrodynamic body forces. The electromagnetic field is
governed by the full Maxwell equations, self-consistently coupled to
the fluid through the space charge and current density. Ion and
electron populations evolve via drift-diffusion-reaction transport.

\begin{definition}[State Vector]
\label{def:state}
The complete state of the coupled system is described by:
\begin{equation*}
  \uu = \bigl(\rho,\; \rho\vv,\; \rho e_t,\;
    n_i,\; n_e,\; \ee,\; \bb\bigr)
  \in \mathbb{R}^{13}
\end{equation*}
where $\rho\;[\mathrm{kg\,m^{-3}}]$ is the gas mass density,
$\vv\;[\mathrm{m\,s^{-1}}]$ the velocity field,
$e_t\;[\mathrm{J\,kg^{-1}}]$ the specific total energy
(internal $+$ kinetic),
$n_i, n_e\;[\mathrm{m^{-3}}]$ the ion and electron number densities,
$\ee\;[\mathrm{V\,m^{-1}}]$ the electric field, and
$\bb\;[\mathrm{T}]$ the magnetic flux density.
\end{definition}

% ---------------------------------------------------------------------------
\subsection{Eqs.~1--4: Coupled Compressible Navier-Stokes and
  Maxwell Equations}
\label{sec:theory:NS_Maxwell}

\subsubsection{Fluid Conservation Laws with EHD Source Terms}

The compressible Navier-Stokes equations, augmented by the EHD body
force $\ff_{\mathrm{EHD}}$ and the Joule heating term $\jj \cdot \ee$,
take the form:

\begin{equation}
\boxed{
  \frac{\partial \rho}{\partial t}
  + \nabla \cdot (\rho \vv) = 0
}
\tag{1}
\label{eq:NS1}
\end{equation}

\noindent
\textbf{Eq.~(1): Mass continuity.} The gas mass density $\rho$ is
conserved; ionization and recombination redistribute particles between
neutral, ion, and electron populations but do not create or destroy
mass at the macroscopic level (the mass of liberated electrons is
negligible, $m_e/m_n \sim 10^{-5}$).

\begin{equation}
\boxed{
  \frac{\partial (\rho \vv)}{\partial t}
  + \nabla \cdot \bigl(\rho \vv \otimes \vv + p\,\mathbf{I}\bigr)
  = \nabla \cdot \bm{\tau} + \ff_{\mathrm{EHD}}
}
\tag{2}
\label{eq:NS2}
\end{equation}

\noindent
\textbf{Eq.~(2): Momentum conservation.} Here $p\;[\mathrm{Pa}]$ is
the thermodynamic pressure (ideal gas: $p = \rho R_{\mathrm{sp}} T$
with specific gas constant $R_{\mathrm{sp}}$ and temperature $T$),
$\mathbf{I}$ the identity tensor, and $\bm{\tau}$ the Newtonian
viscous stress tensor:
\begin{equation}
  \bm{\tau} = \eta\bigl(\nabla\vv + (\nabla\vv)^T\bigr)
  - \frac{2}{3}\eta(\nabla \cdot \vv)\,\mathbf{I}
  \label{eq:viscous_stress}
\end{equation}
with dynamic viscosity $\eta\;[\mathrm{Pa\,s}]$. The EHD body force
comprises three physically distinct contributions:
\begin{equation}
  \ff_{\mathrm{EHD}}
  = \underbrace{q_i n_i \ee}_{\text{ion Coulomb}}
  \;-\; \underbrace{q_e n_e \ee}_{\text{electron Coulomb}}
  \;+\; \underbrace{\frac{1}{2}\nabla\!\left[
    \rho\frac{\partial\varepsilon_r}{\partial\rho}
    \varepsilon_0 |\ee|^2
  \right]}_{\text{electrostriction}}
  \label{eq:f_EHD}
\end{equation}
where $q_i = +e$ (elementary charge) for singly ionized species,
$q_e = -e$, and $\varepsilon_r(\rho)$ is the density-dependent
relative permittivity (Clausius-Mossotti). In atmospheric air at
the field strengths of interest ($|\ee| \sim 10^6\;\mathrm{V/m}$),
the electrostriction term is smaller than the Coulomb terms by
$O(\varepsilon_0 E^2 / p) \sim 10^{-4}$ but is retained for
completeness and because it becomes significant near density gradients.

\begin{equation}
\boxed{
  \frac{\partial (\rho e_t)}{\partial t}
  + \nabla \cdot \bigl[(\rho e_t + p)\vv\bigr]
  = \nabla \cdot (\bm{\tau} \cdot \vv)
  - \nabla \cdot \qq + \jj \cdot \ee
}
\tag{3}
\label{eq:NS3}
\end{equation}

\noindent
\textbf{Eq.~(3): Energy conservation.} The total energy
$\rho e_t = \rho(c_v T + |\vv|^2/2)$ evolves under viscous work,
heat conduction $\qq = -\kappa\nabla T$ (Fourier's law with
thermal conductivity $\kappa\;[\mathrm{W\,m^{-1}\,K^{-1}}]$), and
Ohmic heating $\jj \cdot \ee$. The work done by the EHD body force
$\ff_{\mathrm{EHD}} \cdot \vv$ is implicitly contained in the
left-hand-side convective flux through the pressure and momentum
coupling.

\subsubsection{Maxwell Equations with Self-Consistent Space Charge}

The electromagnetic field is governed by the full Maxwell equations,
not the electrostatic approximation. While the magnetic field effects
are weak in atmospheric EHD ($v/c \sim 10^{-6}$), retaining the full
system is essential for (a)~consistency of the energy balance and
(b)~correct treatment of displacement current effects during
nanosecond-timescale voltage transients in the feedback loop.

\begin{equation}
\boxed{
  \nabla \times \ee = -\frac{\partial \bb}{\partial t}, \quad
  \nabla \times \bb = \mu_0 \jj
    + \mu_0 \varepsilon_0 \frac{\partial \ee}{\partial t}, \quad
  \nabla \cdot \ee = \frac{\rho_q}{\varepsilon_0}, \quad
  \nabla \cdot \bb = 0
}
\tag{4}
\label{eq:Maxwell}
\end{equation}

\noindent
\textbf{Eq.~(4): Maxwell equations.} The net space-charge density is
$\rho_q = q_i n_i + q_e n_e = e(n_i - n_e)$ (for singly ionized
species with $q_e = -e$). The current density comprises drift,
diffusion, and convective contributions:
\begin{equation}
  \jj = e\,n_i(\vv + \mu_i \ee) - e\,D_i\nabla n_i
      - e\,n_e(\vv - \mu_e \ee) + e\,D_e\nabla n_e
  \label{eq:current_density}
\end{equation}
where $\mu_i\;[\mathrm{m^2\,V^{-1}\,s^{-1}}]$ and
$\mu_e\;[\mathrm{m^2\,V^{-1}\,s^{-1}}]$ are the ion and electron
mobilities, and $D_{i,e}$ the corresponding diffusion coefficients.
Diffusion currents are subdominant in the bulk but important in
boundary layers and sheath regions.

\begin{remark}[Displacement Current]
\label{rem:displacement}
During feedback actuation, the electrode voltage changes on a
timescale $\Delta t \sim 1\;\mathrm{ns}$. For a gap $d \sim
1\;\mathrm{cm}$ and $\Delta V \sim 100\;\mathrm{V}$, the
displacement current density
$\varepsilon_0 \partial E/\partial t \sim \varepsilon_0 \Delta V /
(d\,\Delta t) \sim 10^{-1}\;\mathrm{A\,m^{-2}}$ is comparable to
the conduction current
$j_{\mathrm{cond}} \sim e\,n_i\,\mu_i\,E \sim 10^{-1}$--$10^{0}\;
\mathrm{A\,m^{-2}}$. This justifies retaining the full Maxwell
system rather than the electrostatic approximation.
\end{remark}

% ---------------------------------------------------------------------------
\subsection{Eqs.~5--6: Ion and Electron Drift-Diffusion-Reaction
  Transport}
\label{sec:theory:species}

The charged-species number densities evolve according to:

\begin{equation}
\boxed{
  \frac{\partial n_i}{\partial t}
  + \nabla \cdot \bigl(
    n_i \vv + n_i \mu_i \ee - D_i \nabla n_i
  \bigr)
  = \alpha_T |\bm{\Gamma}_e| - \beta_{\mathrm{rec}}\, n_i n_e
}
\tag{5}
\label{eq:species_ion}
\end{equation}

\begin{equation}
\boxed{
  \frac{\partial n_e}{\partial t}
  + \nabla \cdot \bigl(
    n_e \vv - n_e \mu_e \ee - D_e \nabla n_e
  \bigr)
  = \alpha_T |\bm{\Gamma}_e| - \beta_{\mathrm{rec}}\, n_i n_e
}
\tag{6}
\label{eq:species_electron}
\end{equation}

\noindent
\textbf{Eqs.~(5)--(6): Species transport.} Here:
\begin{itemize}
  \item $D_{i,e}\;[\mathrm{m^2\,s^{-1}}]$ are the ion/electron
    diffusion coefficients, related to mobilities by the Einstein
    relation $D_{i,e} = \mu_{i,e} k_B T_{i,e} / e$;
  \item $\alpha_T(|\ee|/N)\;[\mathrm{m^{-1}}]$ is the Townsend
    first ionization coefficient, a strongly nonlinear function of
    the reduced electric field $|\ee|/N$ (where $N$ is the neutral
    number density), typically modeled as
    $\alpha_T = A_T \exp(-B_T N / |\ee|)$ with gas-dependent
    constants $A_T, B_T$ \cite{Raizer1991};
  \item $|\bm{\Gamma}_e| = n_e \mu_e |\ee|$ is the magnitude of the
    electron drift flux;
  \item $\beta_{\mathrm{rec}}\;[\mathrm{m^3\,s^{-1}}]$ is the
    electron-ion recombination coefficient.
\end{itemize}

The ionization source term $S_{\mathrm{ion}} = \alpha_T |\bm{\Gamma}_e|$
embodies the Townsend avalanche mechanism: electrons drifting in the
field gain sufficient energy between collisions to ionize neutral
molecules, creating new ion-electron pairs. This positive-feedback
mechanism is the primary driver of the ionization instability
analyzed in Section~\ref{sec:stability:instabilities}.

\begin{remark}[Closure and Assumptions]
\label{rem:closure}
The system (1)--(6) is closed by the ideal gas equation of state
$p = \rho R_{\mathrm{sp}} T$ and the caloric relation
$e_t = c_v T + |\vv|^2/2$. We assume: (i)~single positive ion
species (e.g., $\mathrm{N}_2^+$ in air); (ii)~no negative ions
(valid for dry air at the timescales of interest,
$\tau \ll 1/\nu_{\mathrm{att}}$ where $\nu_{\mathrm{att}}$ is the
electron attachment frequency); (iii)~local thermodynamic equilibrium
for the neutral gas (the heavy-particle temperature $T$ is
well-defined), while electrons may be non-equilibrium
($T_e \neq T$); (iv)~the gas is optically thin (radiative losses
are negligible).
\end{remark}

% ---------------------------------------------------------------------------
\subsection{Eqs.~7--8: Entropy Production Formalism}
\label{sec:entropy}

The local balance of entropy per unit volume, derived from the
Gibbs relation $T\,ds = de + p\,d(1/\rho) - \sum_\alpha \mu_\alpha
\,dn_\alpha$ applied to the system (1)--(6), yields:

\begin{equation}
\boxed{
  \sdot
  = \underbrace{\frac{\bm{\tau} : \nabla\vv}{T}}_{\sigma_{\mathrm{visc}}}
  + \underbrace{\frac{\kappa\,|\nabla T|^2}{T^2}}_{\sigma_{\mathrm{therm}}}
  + \underbrace{\frac{\jj \cdot \ee
    - \ff_{\mathrm{EHD}} \cdot \vv}{T}}_{\sigma_{\mathrm{Joule\text{-}EHD}}}
  + \underbrace{\frac{W_{\mathrm{ion}}
    S_{\mathrm{ion}}}{T}}_{\sigma_{\mathrm{ion}}}
}
\tag{7}
\label{eq:entropy_local}
\end{equation}

\noindent
\textbf{Eq.~(7): Local entropy production rate.} Each term represents
a distinct irreversible channel:

\begin{enumerate}
  \item \textbf{Viscous dissipation} $\sigma_{\mathrm{visc}}$:
    Conversion of directed kinetic energy to heat via internal
    friction. Scales as $\eta |\nabla\vv|^2 / T$. Dominant in
    shear layers and near electrode surfaces.

  \item \textbf{Thermal conduction} $\sigma_{\mathrm{therm}}$:
    Irreversibility associated with heat flow down temperature
    gradients. Scales as $\kappa |\nabla T|^2 / T^2$. Important
    near the corona sheath where Joule heating creates steep
    temperature gradients.

  \item \textbf{Joule-EHD dissipation} $\sigma_{\mathrm{Joule\text{-}EHD}}$:
    The net irreversible electromagnetic energy deposition. The term
    $\jj \cdot \ee$ represents total electromagnetic energy input;
    subtracting $\ff_{\mathrm{EHD}} \cdot \vv$ (reversible work
    done by EHD forces to accelerate the flow) leaves the
    irreversible Joule component. This is the largest term in
    typical EHD systems.

  \item \textbf{Ionization entropy} $\sigma_{\mathrm{ion}}$:
    Irreversibility of the ionization process itself.
    $W_{\mathrm{ion}}\;[\mathrm{J}]$ is the effective ionization
    energy per event (including excitation losses), and
    $S_{\mathrm{ion}} = \alpha_T |\bm{\Gamma}_e|$ is the ionization
    rate per unit volume. This term is concentrated in the corona
    sheath ($<1\;\mathrm{mm}$ from the emitter).
\end{enumerate}

The global entropy production rate is:

\begin{equation}
\boxed{
  \Sdot = \int_V \sdot \, dV
  = \int_V \bigl(
    \sigma_{\mathrm{visc}}
    + \sigma_{\mathrm{therm}}
    + \sigma_{\mathrm{Joule\text{-}EHD}}
    + \sigma_{\mathrm{ion}}
  \bigr) dV \geq 0
}
\tag{8}
\label{eq:entropy_global}
\end{equation}

\noindent
\textbf{Eq.~(8): Global second law.} The inequality is the Clausius
form of the second law of thermodynamics. The equality holds only
for the (unrealizable) reversible limit. The crucial observation is:
\emph{while the global inequality is inviolable, the spatial
distribution of $\sdot$ within $V$ is not prescribed by thermodynamics
and is therefore controllable by the boundary conditions}---specifically,
by the time-dependent electrode voltage $V_k(t)$.

% ---------------------------------------------------------------------------
\subsection{Proposition 1: The Entropy Redistribution Principle}
\label{sec:theory:theorem1}

We now state and prove the central thermodynamic result.

\begin{definition}[Drift Region]
\label{def:drift_region}
The drift region $\Omega_d \subset V$ is the spatial domain where
ions undergo directed drift from emitter to collector, bounded by
the corona sheath ($\Omega_c$, where ionization is active) and the
collector boundary layer ($\Omega_b$). The total volume decomposes
as $V = \Omega_c \cup \Omega_d \cup \Omega_b \cup \Omega_{\mathrm{ext}}$.
\end{definition}

\begin{definition}[EHD Propulsive Efficiency]
\label{def:efficiency}
The propulsive efficiency is:
\begin{equation}
  \eta_{\mathrm{EHD}}
  \equiv \frac{F_{\mathrm{thrust}} \cdot v_{\infty}}{P_{\mathrm{elec}}}
  = \frac{F_{\mathrm{thrust}} \cdot v_{\infty}}
         {T\,\Sdot + F_{\mathrm{thrust}} \cdot v_{\infty}}
  \label{eq:efficiency}
\end{equation}
where $F_{\mathrm{thrust}}$ is the net body-integrated EHD force,
$v_{\infty}$ is the freestream velocity, and $P_{\mathrm{elec}}$
is the total electrical power input. Eq.~(\ref{eq:efficiency})
follows from the first law, partitioning $P_{\mathrm{elec}}$ into
useful thrust power and irreversible entropy production. For static
thrust conditions ($v_{\infty} = 0$), which are typical of
laboratory EHD devices, this definition is degenerate; in that case,
the relevant figure of merit is the thrust-to-power ratio
$F_{\mathrm{thrust}} / P_{\mathrm{elec}}\;[\mathrm{N/W}]$.
\end{definition}

The efficiency $\eta_{\mathrm{EHD}}$ is maximized by minimizing
$\Sdot$---but $\Sdot$ cannot be made zero. It \emph{can}, however,
be redistributed.

\begin{proposition}[Entropy Redistribution Principle]
\label{thm:entropy_redistribution}
Consider the coupled system \textnormal{(1)--(6)} on a domain $V$
with time-dependent Dirichlet boundary conditions on the electrode
surfaces:
\begin{equation}
  \ee\big|_{\partial V_k}
  = \ee_0^{(k)} + \delta\ee^{(k)}(t), \quad k = 1,\ldots,m
  \label{eq:BC}
\end{equation}
where $\ee_0^{(k)}$ is the steady-state field and
$\delta\ee^{(k)}(t)$ is the control perturbation applied to the
$k$-th electrode segment. Let $\gamma_{\max}$ be the maximum
growth rate of linear instabilities of the uncontrolled steady state
$\uu_0$. If the control signal has bandwidth
$f_{\mathrm{ctrl}} > \gamma_{\max}/(2\pi)$ and satisfies
\begin{equation}
  \|\delta\ee\|_{\infty}
  \leq C_{\mathrm{stab}}\,\gamma_{\max}^{-1}\,
  \max_{\mathbf{k}} |\hat{n}_i'(\mathbf{k}, t=0)|
  \label{eq:control_bound}
\end{equation}
for a constant $C_{\mathrm{stab}}$ depending on the electrode
geometry, then:

\noindent\textnormal{(i)} The perturbation $\uu' = \uu - \uu_0$
satisfies $\|\uu'(t)\| \to 0$ as $t \to \infty$ (asymptotic
stability);

\noindent\textnormal{(ii)} The entropy production in the drift
region satisfies
\begin{equation}
  \int_{\Omega_d} \sdot^{\,\mathrm{ctrl}} \, dV
  \leq \int_{\Omega_d} \sdot^{\,\mathrm{lam}} \, dV
  + O\!\left(\frac{\gamma_{\max}}{f_{\mathrm{ctrl}}}\right)^2
  \label{eq:entropy_drift_bound}
\end{equation}
where $\sdot^{\,\mathrm{lam}}$ is the entropy production of the
(unstable) laminar base state, which satisfies
\begin{equation}
  \int_{\Omega_d} \sdot^{\,\mathrm{lam}} \, dV
  \ll \int_{\Omega_d} \sdot^{\,\mathrm{turb}} \, dV
  \label{eq:lam_vs_turb}
\end{equation}
with $\sdot^{\,\mathrm{turb}}$ the turbulent entropy production rate;

\noindent\textnormal{(iii)} The global constraint $\Sdot \geq 0$
is preserved, with the entropy cost of control satisfying
\begin{equation}
  \Delta\Sdot^{\,\mathrm{ctrl}}
  \equiv \Sdot^{\,\mathrm{ctrl}} - \Sdot^{\,\mathrm{lam}}
  = \int_{\partial V_k} \frac{|\delta\jj|^2}{\sigma_{\mathrm{eff}}\,T}
    \, dA + O(\|\delta\ee\|^3)
  \label{eq:entropy_cost}
\end{equation}
which is localized at the electrode surfaces and is negligible
compared to the entropy \emph{saved} in $\Omega_d$.
\end{proposition}

\textit{Argument.}
We decompose the entropy production into controlled and uncontrolled
components using $\uu = \uu_0 + \uu'$ and $\ee = \ee_0 + \delta\ee$.

\textit{Part (i):} The linearized dynamics satisfy
$\partial_t \uu' = \Lop\,\uu' + \mathbf{B}\,\delta\ee$
(cf.\ Eq.~\ref{eq:linearized_full} below). The control law
$\delta\ee = -\kk\,\mathbf{C}\,\uu'$ (Section~\ref{sec:stability:control})
yields $\partial_t \uu' = (\Lop - \mathbf{B}\kk\mathbf{C})\,\uu'$.
By Proposition~\ref{thm:controllability} below, the gain matrix $\kk$
can be chosen such that all eigenvalues of $\Lop - \mathbf{B}\kk\mathbf{C}$
have negative real part, ensuring asymptotic stability of the
\emph{finite-dimensional truncation}.

\begin{remark}[Truncation caveat]
This argument applies to the Galerkin truncation retaining $N_m$
modes. Spillover instability---where controlling the resolved modes
destabilizes unresolved ones---is a known risk in distributed-parameter
systems and must be assessed numerically for any specific electrode
configuration. A rigorous infinite-dimensional treatment would require
proving well-posedness and establishing $H^1$-norm bounds on the
perturbation, which is beyond the scope of this work.
\end{remark}

\textit{Part (ii):} In the stabilized state, $\uu' \to 0$ and the
flow approaches the laminar base state $\uu_0$. The entropy production
in $\Omega_d$ approaches $\sdot^{\,\mathrm{lam}}$ up to corrections
of order $\|\uu'\|^2 \sim (\gamma_{\max}/f_{\mathrm{ctrl}})^2$
from residual perturbations within the control bandwidth.
The inequality (\ref{eq:lam_vs_turb}) is motivated by the standard
result that turbulent viscous dissipation exceeds laminar dissipation
by a factor proportional to $\mathrm{Re}^{1/2}$ to $\mathrm{Re}^{3/4}$
in Newtonian channel flows \cite{Pope2000}. We note that direct
transfer of this scaling to EHD systems with space-charge body forces
is an assumption that requires validation through direct numerical
simulation of the coupled system.

\textit{Part (iii):} The additional Joule dissipation from the
control perturbation is $\delta\jj \cdot \delta\ee \approx
|\delta\jj|^2 / \sigma_{\mathrm{eff}}$, localized at the electrode
surfaces where $\delta\ee$ is applied. The assumption
$\|\delta\ee\| / \|\ee_0\| \sim 10^{-2}$--$10^{-1}$ is an estimate
based on the linearized control gain; whether this amplitude is
achievable depends on the specific system parameters and must be
verified for each configuration. Under this assumption, the control
contributes $\sim 10^{-4}$--$10^{-2}$ of the total power, hence
the entropy cost is small. The second law is preserved because
$\Sdot^{\,\mathrm{ctrl}} = \Sdot^{\,\mathrm{lam}} +
\Delta\Sdot^{\,\mathrm{ctrl}} > 0$.

\begin{remark}[Physical Interpretation]
Proposition~\ref{thm:entropy_redistribution} argues that feedback
control does not \emph{reduce} global entropy production (which would
violate the second law) but \emph{redirects} it: entropy that would
have been produced in $\Omega_d$ (destroying coherent thrust) is
instead produced at the electrode boundaries (as a small, localized
Joule cost of the control action).
\end{remark}

% ============================================================================
%  3. STABILITY ANALYSIS
% ============================================================================
\section{Stability Analysis}
\label{sec:stability}

% ---------------------------------------------------------------------------
\subsection{Linearization of the Coupled System}
\label{sec:stability:linearization}

Let $\uu_0(\mathbf{x})$ be a steady-state solution of (1)--(6)
representing the laminar ion drift between asymmetric electrodes
under constant applied voltage $V_0$. We write
$\uu = \uu_0 + \epsilon\,\uu'$ with $\epsilon \ll 1$ and retain
$O(\epsilon)$ terms:

\begin{equation}
  \frac{\partial \uu'}{\partial t} = \Lop[\uu_0]\,\uu'
  \label{eq:linearized_full}
\end{equation}

where $\Lop$ is the linearized operator, a partial differential
operator in the spatial variables whose coefficients depend on the
base state $\uu_0$. For modal analysis, we decompose:

\begin{equation}
  \uu'(\mathbf{x}, t)
  = \sum_{\mathbf{k}} \hat{\uu}(\mathbf{k})\,
    e^{i\mathbf{k}\cdot\mathbf{x} + \lambda(\mathbf{k})\,t}
  \label{eq:modal_decomposition}
\end{equation}

yielding the eigenvalue problem:

\begin{equation}
  \lambda(\mathbf{k})\,\hat{\uu}
  = \hat{\Lop}(\mathbf{k})\,\hat{\uu}
  \label{eq:dispersion_EVP}
\end{equation}

The system is linearly unstable if $\exists\,\mathbf{k}$ such that
$\mathrm{Re}[\lambda(\mathbf{k})] > 0$.

% ---------------------------------------------------------------------------
\subsection{Three Dominant Instability Mechanisms}
\label{sec:stability:instabilities}

The spectrum $\{\lambda(\mathbf{k})\}$ of the coupled system
exhibits three families of unstable modes, each traceable to a
distinct physical mechanism. We derive the growth rate for each
from the governing equations.

\subsubsection{(I) Electrohydrodynamic (Coulomb-Driven) Instability}

Consider a perturbation $\delta n_i > 0$ in the ion density within
the drift region. By Gauss's law (Eq.~4),
$\nabla \cdot \delta\ee = e\,\delta n_i / \varepsilon_0$, producing
a perturbed field $|\delta\ee| \sim e\,\delta n_i\,L / \varepsilon_0$
over the scale $L$. This field exerts a Coulomb force
$\delta f \sim e\,n_{i,0}\,\delta E$ on the background ion density,
driving a momentum perturbation. The growth rate follows from
balancing the acceleration $\rho_0 \partial v' / \partial t
\sim \rho_0 \gamma_{\mathrm{EHD}} v'$ against the Coulomb force
$e\,n_{i,0}\,\delta E \sim e^2 n_{i,0} \delta n_i L /
\varepsilon_0$:

\begin{equation}
  \gamma_{\mathrm{EHD}}
  = \left(\frac{e^2 n_{i,0}^2}{\varepsilon_0 \rho_0}\right)^{1/2}
  \equiv \omega_{\mathrm{pi}}^{\mathrm{eff}}
  \label{eq:gamma_EHD}
\end{equation}

where we identify $\omega_{\mathrm{pi}}^{\mathrm{eff}}$ as an
effective ion plasma frequency for the weakly ionized gas. We note
that this estimate neglects ion-neutral collision damping
($\nu_{\mathrm{in}} \sim 10^{9}$--$10^{10}\;\mathrm{s^{-1}}$ at
atmospheric pressure), which can reduce the effective growth rate
by an order of magnitude or more. For atmospheric air with
$n_{i,0} \sim 10^{18}\;\mathrm{m^{-3}}$ and
$\rho_0 \approx 1.2\;\mathrm{kg\,m^{-3}}$:

\begin{equation}
  \gamma_{\mathrm{EHD}}
  = \left(
    \frac{(1.6 \times 10^{-19})^2 \times (10^{18})^2}
         {8.85 \times 10^{-12} \times 1.2}
  \right)^{1/2}
  \approx 4.9 \times 10^7\;\mathrm{s^{-1}}
  \label{eq:gamma_EHD_num}
\end{equation}

\subsubsection{(II) Ionization-Driven (Townsend) Instability}

A local charge density perturbation $\delta n_i$ enhances the
electric field in the upstream region (toward the emitter), which
exponentially increases the ionization rate through the Townsend
coefficient $\alpha_T \propto \exp(-B_T N / |\ee|)$. Linearizing
the ionization source term in Eq.~(5):

\begin{equation}
  \delta S_{\mathrm{ion}}
  = \frac{\partial S_{\mathrm{ion}}}{\partial |\ee|}\,\delta|\ee|
  + \frac{\partial S_{\mathrm{ion}}}{\partial n_e}\,\delta n_e
  \label{eq:ionization_perturbation}
\end{equation}

The first term provides positive feedback: increased $\delta|\ee|$
$\to$ increased ionization $\to$ increased $\delta n_i, \delta n_e$
$\to$ further increased $\delta|\ee|$. The growth rate is dominated
by the field-sensitivity of $\alpha_T$:

\begin{equation}
  \gamma_{\mathrm{ion}}
  = \mu_e E_0 \cdot \alpha_T(E_0)
    \cdot \frac{B_T N}{E_0}
  \label{eq:gamma_ion}
\end{equation}

The factor $B_T N / E_0$ is the logarithmic sensitivity of the
Townsend coefficient to field variations. For air at STP with
$E_0 \sim 3 \times 10^6\;\mathrm{V/m}$ (near breakdown),
$\mu_e \approx 0.03\;\mathrm{m^2\,V^{-1}\,s^{-1}}$,
$\alpha_T \approx 2 \times 10^4\;\mathrm{m^{-1}}$, and
$B_T N / E_0 \approx 5$:

\begin{equation}
  \gamma_{\mathrm{ion}}
  \approx 0.03 \times 3 \times 10^6 \times 2 \times 10^4 \times 5
  = 9.0 \times 10^9\;\mathrm{s^{-1}}
  \label{eq:gamma_ion_num}
\end{equation}

\begin{remark}[Limitations of the growth rate estimate]
Equation~(\ref{eq:gamma_ion}) is a scaling estimate, not a solution
of the full dispersion relation. It neglects: (i)~electron diffusion
damping ($D_e k^2$), which stabilizes short-wavelength modes;
(ii)~recombination damping ($\beta_{\mathrm{rec}} n_{i,0}$); and
(iii)~the fact that at $E_0 \sim 3 \times 10^6\;\mathrm{V/m}$
(near breakdown), the Townsend avalanche model is at the edge of its
validity---streamer formation may dominate at these field strengths
\cite{Raizer1991}. A proper dispersion analysis solving the
linearized species-Poisson system would yield a wavenumber-dependent
growth rate $\gamma_{\mathrm{ion}}(k)$ with a maximum at finite $k$,
likely giving a lower effective $\gamma_{\max}$. The value
$9.0 \times 10^9\;\mathrm{s^{-1}}$ should therefore be interpreted
as an \emph{upper bound estimate}.
\end{remark}

\subsubsection{(III) Kelvin-Helmholtz (Shear) Instability}

The ion wind jet emanating from the corona region creates a velocity
shear layer at the jet boundary. The classical Kelvin-Helmholtz
growth rate for an inviscid shear layer with velocity difference
$\Delta v$ and shear layer thickness $\delta_s$ is:

\begin{equation}
  \gamma_{\mathrm{KH}}
  = \frac{k\,\Delta v}{2}
  \quad \text{for } k\,\delta_s \lesssim 1
  \label{eq:gamma_KH}
\end{equation}

where $k$ is the perturbation wavenumber. We note that this inviscid,
incompressible estimate becomes inaccurate at high Mach numbers
($M > 0.5$), where compressibility strongly stabilizes KH modes.
For typical EHD \emph{bulk gas} velocities $\Delta v \sim
1$--$10\;\mathrm{m/s}$ (as measured in experiments \cite{Xu2018,
Moreau2007}; note that the ion drift velocity $\mu_i E_0 \sim
600\;\mathrm{m/s}$ is that of individual ions, not the bulk flow)
and $\delta_s \sim 1\;\mathrm{mm}$:

\begin{equation}
  \gamma_{\mathrm{KH}}
  \sim \frac{1\text{--}10}{2 \times 10^{-3}}
  = 5 \times 10^2\text{--}5 \times 10^3\;\mathrm{s^{-1}}
  \label{eq:gamma_KH_num}
\end{equation}

\subsubsection{Dominant Growth Rate}

Comparing the three mechanisms:

\begin{equation}
  \gamma_{\mathrm{ion}} \;\lesssim\; 10^{9}\;\mathrm{s^{-1}}
  \;\gg\; \gamma_{\mathrm{EHD}} \;\lesssim\; 5 \times 10^{7}\;\mathrm{s^{-1}}
  \;\gg\; \gamma_{\mathrm{KH}} \;\sim\; 10^{3}\text{--}10^{4}\;\mathrm{s^{-1}}
  \label{eq:gamma_hierarchy}
\end{equation}

The ionization-driven instability is the fastest and therefore the
rate-limiting process for control:

\begin{equation}
\boxed{
  \gamma_{\max}
  = \gamma_{\mathrm{ion}}
  \sim 10^{9}\;\mathrm{s^{-1}}
}
\label{eq:gamma_max}
\end{equation}

% ---------------------------------------------------------------------------
\subsection{Mathematical Necessity of GHz Sampling}
\label{sec:stability:nyquist}

The feedback control system must satisfy two fundamental constraints:

\begin{proposition}[Sampling Rate Bound]
\label{prop:sampling}
For a linear feedback controller to stabilize all unstable modes
of the linearized system \textnormal{(\ref{eq:linearized_full})},
the sampling frequency $f_s$ must satisfy:

\noindent\textnormal{(i)} \textbf{Nyquist condition}
(mode observability):
\begin{equation}
  f_s > 2\,f_{\max}
  = \frac{2\,\gamma_{\max}}{2\pi}
  = \frac{\gamma_{\max}}{\pi}
  \approx \frac{10^9}{\pi}
  \approx 3.2 \times 10^8\;\mathrm{Hz}
  \label{eq:nyquist}
\end{equation}

\noindent\textnormal{(ii)} \textbf{Control authority condition}
(mode suppression):
\begin{equation}
  f_s > \frac{\gamma_{\max}}{-\ln(1 - 2\gamma_{\max}/f_s)}
  \;\;\Rightarrow\;\;
  f_s \gtrsim 10\,\gamma_{\max}/(2\pi)
  \approx 1.6 \times 10^9\;\mathrm{Hz}
  \label{eq:control_authority}
\end{equation}
where the factor of $\sim 10$ arises from requiring the closed-loop
eigenvalue magnitude to be reduced by at least $e^{-1}$ per sampling
period.
\end{proposition}

\begin{proof}
Part (i) is the Shannon-Nyquist theorem applied to the observation
of the most rapidly oscillating/growing mode. Part (ii) follows from
discrete-time control theory: the sampled system
$\uu'_{n+1} = e^{\Lop/f_s}\uu'_n + \mathbf{B}_d\,\delta\ee_n$
requires $\|e^{\Lop/f_s}\| < 1$ after feedback, which for the
dominant eigenvalue $\gamma_{\max}$ gives
$e^{\gamma_{\max}/f_s}(1 - K_{\mathrm{eff}}) < 1$. For
$K_{\mathrm{eff}} \leq 1$ (linear regime), this yields
$f_s > \gamma_{\max} / \ln(1/(1-K_{\mathrm{eff}}))$. Requiring
$K_{\mathrm{eff}} < 0.8$ for robustness gives
$f_s > \gamma_{\max}/\ln(5) \approx 0.62\,\gamma_{\max}$, or
$f_s \gtrsim 6.2 \times 10^8\;\mathrm{Hz}$. With engineering margin
and accounting for pipeline latency, we require
$f_s \geq 1\;\mathrm{GHz}$.
\end{proof}

\begin{remark}
This is a rigorous lower bound, not an empirical choice. No
sub-GHz control system can stabilize the ionization-driven instability
in atmospheric-pressure EHD devices. This mathematical fact is the
primary motivation for the FPGA architecture of
Section~\ref{sec:fpga}.
\end{remark}

% ---------------------------------------------------------------------------
\subsection{Proposition 2: Controllability of the Coupled EHD System}
\label{sec:stability:control}

We now address the fundamental question: given an FPGA controller
with sufficient bandwidth, \emph{can} the unstable modes be stabilized?

Consider the linearized system in the form:
\begin{align}
  \frac{d\uu'}{dt} &= \Lop\,\uu' + \mathbf{B}\,\delta\bm{u}(t)
  \label{eq:state_eq} \\
  \bm{y}(t) &= \mathbf{C}\,\uu'
  \label{eq:output_eq}
\end{align}
where $\uu' \in \mathbb{R}^n$ is the (truncated) state vector,
$\delta\bm{u}(t) \in \mathbb{R}^m$ is the control input
(voltage perturbations on $m$ electrode segments),
$\bm{y}(t) \in \mathbb{R}^p$ is the measurement output
($p$ current/field sensors), $\mathbf{B} \in \mathbb{R}^{n \times m}$
is the input (actuator) matrix, and
$\mathbf{C} \in \mathbb{R}^{p \times n}$ is the output (sensor)
matrix.

\begin{definition}[Controllability and Observability Matrices]
\label{def:kalman}
The controllability and observability matrices are:
\begin{align}
  \mathcal{C} &= \bigl[\mathbf{B},\; \Lop\mathbf{B},\;
    \Lop^2\mathbf{B},\; \ldots,\; \Lop^{n-1}\mathbf{B}\bigr]
    \in \mathbb{R}^{n \times mn}
    \label{eq:controllability_matrix} \\
  \mathcal{O} &= \bigl[\mathbf{C}^T,\;
    (\Lop\mathbf{C})^T,\; (\Lop^2\mathbf{C})^T,\;
    \ldots,\; (\Lop^{n-1}\mathbf{C})^T\bigr]^T
    \in \mathbb{R}^{pn \times n}
    \label{eq:observability_matrix}
\end{align}
\end{definition}

\begin{proposition}[EHD Controllability---Kalman Rank Condition]
\label{thm:controllability}
The linearized coupled Navier-Stokes--Maxwell system
\textnormal{(\ref{eq:state_eq})--(\ref{eq:output_eq})} with $m$
independently driven electrode segments and $p$ current/field sensors
is output-feedback stabilizable if and only if:

\noindent\textnormal{(i)} \textbf{Controllability:}
\begin{equation}
  \mathrm{rank}\,\mathcal{C}
  = \mathrm{rank}\bigl[\mathbf{B},\; \Lop\mathbf{B},\;
    \ldots,\; \Lop^{n-1}\mathbf{B}\bigr] = n
  \label{eq:controllability_rank}
\end{equation}

\noindent\textnormal{(ii)} \textbf{Observability:}
\begin{equation}
  \mathrm{rank}\,\mathcal{O}
  = \mathrm{rank}\bigl[\mathbf{C}^T,\;
    (\Lop\mathbf{C})^T,\; \ldots,\;
    (\Lop^{n-1}\mathbf{C})^T\bigr]^T = n
  \label{eq:observability_rank}
\end{equation}

\noindent
Equivalently, by the Hautus lemma, the system is stabilizable
iff for every unstable eigenvalue $\lambda_k$ with
$\mathrm{Re}(\lambda_k) > 0$:
\begin{equation}
  \mathrm{rank}\begin{bmatrix}
    \lambda_k \mathbf{I} - \Lop & \mathbf{B}
  \end{bmatrix} = n
  \quad \text{and} \quad
  \mathrm{rank}\begin{bmatrix}
    \lambda_k \mathbf{I} - \Lop \\ \mathbf{C}
  \end{bmatrix} = n
  \label{eq:hautus}
\end{equation}
\end{proposition}

\textit{Justification.}
The conditions are an application of the Kalman
controllability/observability decomposition \cite{Kalman1960} to the
finite-dimensional truncation of the linearized PDE system. The
truncation is justified by the spectral gap: only finitely many modes
($N_u$) are unstable ($\mathrm{Re}(\lambda_k) > 0$), while the
remaining modes are exponentially damped by viscosity and diffusion.
We emphasize that this proposition identifies \emph{structural
conditions} for stabilizability; \emph{verification} that these
conditions hold for a specific electrode geometry requires numerical
computation of the matrices $\mathbf{B}$ and $\mathbf{C}$ and
evaluation of the rank conditions, which is deferred to future work.

\textit{Controllability:}
The input matrix $\mathbf{B}$ maps the $m$ electrode voltage
perturbations to field perturbations in the domain. The $j$-th
column of $\mathbf{B}$ is the response of $\uu'$ to a unit voltage
impulse on electrode $j$, which---by linearity of the Poisson
equation in the electrostatic limit---spans a distinct spatial mode
for each geometrically distinct electrode segment. If $m \geq N_u$
and the electrode segments are not geometrically degenerate (i.e.,
they produce linearly independent field perturbations at the locations
of the unstable modes), then $\mathcal{C}$ has rank $\geq N_u$, and
the unstable subspace is controllable.

\textit{Observability:}
The output matrix $\mathbf{C}$ maps the state to measurable currents
and field values at $p$ sensor locations. Each unstable mode produces
a characteristic current signature at the sensors. If $p \geq N_u$
and the sensors are not co-located with nodes of all unstable
eigenmodes, then $\mathcal{O}$ has rank $\geq N_u$.

\textit{Stabilizability:}
Given full-rank $\mathcal{C}$ and $\mathcal{O}$, the separation
principle guarantees that a state estimator (Kalman filter) and
state-feedback controller can be designed independently, with the
combined system achieving any desired closed-loop pole placement
\cite{Anderson1990}. The gain matrix $\kk$ is computed via the
algebraic Riccati equation:
\begin{equation}
  \Lop^T \mathbf{P} + \mathbf{P}\Lop
  - \mathbf{P}\mathbf{B}\mathbf{R}^{-1}\mathbf{B}^T\mathbf{P}
  + \mathbf{Q} = 0
  \label{eq:riccati}
\end{equation}
with $\kk = \mathbf{R}^{-1}\mathbf{B}^T\mathbf{P}$, where
$\mathbf{Q} \succ 0$ and $\mathbf{R} \succ 0$ are the state and
control cost matrices, respectively.

\begin{corollary}[Minimum Electrode Segmentation]
\label{cor:min_electrodes}
If the system has $N_u$ unstable modes, the minimum number of
independently controllable electrode segments required for
stabilizability is $m \geq N_u$. For atmospheric EHD with
$N_u \sim 3$--$5$ dominant unstable modes,
$m \geq 4$ is sufficient.
\end{corollary}

% ---------------------------------------------------------------------------
\subsection{Entropy Production Under Feedback Control:
  The Information-Thermodynamic Bridge}
\label{sec:stability:entropy_ctrl}

Under successful stabilization, the total entropy production rate
decomposes as:

\begin{equation}
  \Sdot^{\,\mathrm{ctrl}}
  = \Sdot^{\,\mathrm{lam}} + \Delta\Sdot^{\,\mathrm{ctrl}}
  \label{eq:entropy_decomp}
\end{equation}

The control cost $\Delta\Sdot^{\,\mathrm{ctrl}}$ has two components:

\begin{enumerate}
  \item \textbf{Electromagnetic cost}: The Joule dissipation of the
    control currents at the electrodes:
    \begin{equation}
      \Delta\Sdot^{\,\mathrm{EM}}
      = \int_V \frac{|\delta\jj|^2}{\sigma_{\mathrm{eff}}\,T}\,dV
      \sim \varepsilon_0 E_0 (\delta E)^2 A\,f_s / T
      \label{eq:entropy_EM}
    \end{equation}
    For $\delta E / E_0 \sim 10^{-2}$, this is $\sim 10^{-4}$ of
    the total power.

  \item \textbf{Information-theoretic (Landauer) cost}: The
    thermodynamic cost of information acquisition and processing:
    \begin{equation}
      \Delta\Sdot^{\,\mathrm{Landauer}}
      = f_s \cdot N_{\mathrm{ch}} \cdot N_{\mathrm{bits}}
        \cdot k_B \ln 2
      \label{eq:entropy_Landauer}
    \end{equation}
    where $N_{\mathrm{ch}}$ is the number of sensor channels. For
    $f_s = 1\;\mathrm{GHz}$, $N_{\mathrm{ch}} = 8$,
    $N_{\mathrm{bits}} = 12$:
    \begin{equation}
      \Delta\Sdot^{\,\mathrm{Landauer}}
      = 10^9 \times 8 \times 12 \times 1.38 \times 10^{-23}
        \times 0.693
      \approx 9.2 \times 10^{-13}\;\mathrm{W/K}
      \label{eq:landauer_numerical}
    \end{equation}
\end{enumerate}

The entropy saved by preventing turbulent transition in the drift
region is:

\begin{equation}
  \Delta\Sdot^{\,\mathrm{saved}}
  = \int_{\Omega_d} \bigl(\sdot^{\,\mathrm{turb}}
    - \sdot^{\,\mathrm{lam}}\bigr)\,dV
  \sim \frac{P_{\mathrm{elec}}}{T}
    \cdot \frac{\Reehd^{1/2} - 1}{\Reehd^{1/2}}
  \label{eq:entropy_saved}
\end{equation}

For a $10\;\mathrm{W}$ device at $T = 300\;\mathrm{K}$ with
$\Reehd \sim 10^4$:

\begin{equation}
  \Delta\Sdot^{\,\mathrm{saved}}
  \sim \frac{10}{300} \times 0.99
  \approx 3.3 \times 10^{-2}\;\mathrm{W/K}
  \label{eq:entropy_saved_num}
\end{equation}

\begin{definition}[Thermodynamic Leverage Factor]
\label{def:leverage}
The thermodynamic leverage of the control system is:
\begin{equation}
  \Lambda
  \equiv \frac{\Delta\Sdot^{\,\mathrm{saved}}}
              {\Delta\Sdot^{\,\mathrm{Landauer}}}
  \approx \frac{3.3 \times 10^{-2}}{9.2 \times 10^{-13}}
  \approx 3.6 \times 10^{10}
  \label{eq:leverage}
\end{equation}
\end{definition}

This factor of $\sim 10^{10}$ quantifies the \emph{theoretical}
thermodynamic efficiency at the Landauer floor: each bit of
information, at the minimum thermodynamic cost, prevents the
dissipation of $\sim 10^{10}$ times that cost in macroscopic entropy.

\begin{remark}[Practical vs.\ theoretical leverage]
The Landauer bound represents the absolute thermodynamic minimum for
information erasure. A real FPGA operating at GHz clock rates
dissipates $P_{\mathrm{FPGA}} \sim 30$--$100\;\mathrm{W}$, giving a
practical leverage factor:
\begin{equation}
  \Lambda_{\mathrm{real}}
  = \frac{\Delta\Sdot^{\,\mathrm{saved}}}{P_{\mathrm{FPGA}}/T}
  \sim \frac{3.3 \times 10^{-2}}{50/300}
  \approx 0.2
  \label{eq:leverage_real}
\end{equation}
This is less than unity for a $10\;\mathrm{W}$ EHD device, meaning
the FPGA itself dissipates more entropy than it saves in the plasma.
The practical viability of the scheme therefore depends on scaling to
higher-power EHD systems ($P_{\mathrm{elec}} \gg P_{\mathrm{FPGA}}$)
where $\Lambda_{\mathrm{real}} > 1$. The Landauer-based leverage
$\Lambda \sim 10^{10}$ quantifies the \emph{fundamental} information-
thermodynamic amplification available in the far-from-equilibrium
plasma, not the practical efficiency of current hardware.
\end{remark}

% ============================================================================
%  4. FPGA CONTROL ARCHITECTURE
% ============================================================================
\section{FPGA Control Architecture}
\label{sec:fpga}

% ---------------------------------------------------------------------------
\subsection{System Requirements Derived from Theory}

The theoretical analysis of Sections~\ref{sec:stability:nyquist}
and \ref{sec:stability:control} imposes precise hardware requirements:

\begin{table}[h]
\centering
\caption{Control system requirements derived from first principles.}
\label{tab:requirements}
\begin{tabular}{@{}llp{3.8cm}@{}}
\toprule
\textbf{Parameter} & \textbf{Requirement} & \textbf{Derivation} \\
\midrule
Sampling rate $f_s$
  & $\geq 1\;\mathrm{GHz}$
  & Prop.~\ref{prop:sampling}: $f_s > 10\gamma_{\max}/(2\pi)$ \\
Feedback latency $\tau_{\mathrm{fb}}$
  & $\leq 10\;\mathrm{ns}$
  & $\tau_{\mathrm{fb}} < \gamma_{\max}^{-1} = 1\;\mathrm{ns}
    \times O(10)$ for phase margin \\
ADC resolution
  & $\geq 12\;\mathrm{bit}$
  & Mode discrimination SNR: $20\log_{10}(2^{12}) = 72\;\mathrm{dB}$ \\
DAC resolution
  & $\geq 14\;\mathrm{bit}$
  & $\delta E/E_0 \sim 10^{-2}$ requires $> 10^{4}$ levels \\
Sensor channels $N_{\mathrm{ch}}$
  & $\geq 2 N_u$
  & Observability: Cor.~\ref{cor:min_electrodes} \\
Actuator channels $m$
  & $\geq N_u$
  & Controllability: Thm.~\ref{thm:controllability} \\
\bottomrule
\end{tabular}
\end{table}

% ---------------------------------------------------------------------------
\subsection{Pipeline Architecture: 7--9 Clock Cycle Latency}
\label{sec:fpga:pipeline}

The feedback loop is implemented as a deeply pipelined datapath on
the FPGA fabric. The logical justification for each stage and its
cycle budget:

\begin{algorithm}[h]
\caption{FPGA Real-Time Feedback Pipeline}
\label{alg:pipeline}
\begin{algorithmic}[1]
\Statex \textbf{Stage 1: Acquisition} [1 cycle]
\State Latch $N_{\mathrm{ch}}$ ADC samples $\{s_k[n]\}_{k=1}^{N_{\mathrm{ch}}}$
\Statex \textit{Justification: RF-ADCs on modern FPGAs (Xilinx Versal,}
\Statex \textit{Intel Agilex) provide synchronous digital output at}
\Statex \textit{$\geq 4\;\mathrm{GS/s}$ with 1-cycle latency.}

\Statex
\Statex \textbf{Stage 2: Baseline Subtraction \& Filtering} [1 cycle]
\State $\delta s_k[n] \gets s_k[n] - \bar{s}_k$ \Comment{Remove DC}
\State Apply single-stage IIR: $\delta\tilde{s}_k \gets \alpha\,\delta s_k[n] + (1-\alpha)\,\delta\tilde{s}_k^{\mathrm{prev}}$
\Statex \textit{Justification: Subtract precomputed steady-state values}
\Statex \textit{(stored in BRAM). IIR filter for noise rejection uses}
\Statex \textit{one DSP slice per channel.}

\Statex
\Statex \textbf{Stage 3: State Estimation (Kalman Predict)} [2 cycles]
\State $\hat{\uu}'_{\mathrm{pred}} \gets \mathbf{A}_d\,\hat{\uu}'_{\mathrm{prev}} + \mathbf{B}_d\,\delta\bm{u}_{\mathrm{prev}}$
\State $\mathbf{P}_{\mathrm{pred}} \gets \mathbf{A}_d\,\mathbf{P}_{\mathrm{prev}}\,\mathbf{A}_d^T + \mathbf{Q}_d$
\Statex \textit{Justification: Discretized state transition ($\mathbf{A}_d = e^{\Lop/f_s}$)}
\Statex \textit{precomputed offline. Matrix dimensions: $N_m \times N_m$}
\Statex \textit{with $N_m \leq 10$. Two cycles for the two matrix-vector}
\Statex \textit{products using pipelined DSP arrays.}

\Statex
\Statex \textbf{Stage 4: Kalman Update \& Mode Projection} [2 cycles]
\State Innovation: $\bm{e} \gets \delta\tilde{\bm{s}} - \mathbf{C}_d\,\hat{\uu}'_{\mathrm{pred}}$
\State Kalman gain: $\mathbf{K}_f \gets \mathbf{P}_{\mathrm{pred}}\,\mathbf{C}_d^T\,(\mathbf{C}_d\,\mathbf{P}_{\mathrm{pred}}\,\mathbf{C}_d^T + \mathbf{R}_d)^{-1}$
\State $\hat{\uu}' \gets \hat{\uu}'_{\mathrm{pred}} + \mathbf{K}_f\,\bm{e}$
\State Mode amplitudes: $a_j \gets [\hat{\uu}']_j$ for $j = 1,\ldots,N_m$
\Statex \textit{Justification: The matrix inverse is precomputed for the}
\Statex \textit{slowly varying $\mathbf{P}$; only the product $\mathbf{K}_f\,\bm{e}$}
\Statex \textit{is computed in real time ($N_m \times p$ MAC operations).}

\Statex
\Statex \textbf{Stage 5: Control Law Computation} [1--2 cycles]
\State $\delta\bm{u}[n] \gets -\kk\,\hat{\uu}'$ \Comment{$m \times N_m$ MAC}
\State Apply saturation: $|\delta u_k| \leq \delta u_{\max}$
\Statex \textit{Justification: Linear state feedback. The gain matrix $\kk$}
\Statex \textit{is loaded from registers (updated by outer loop).}
\Statex \textit{Saturation prevents actuator overload.}

\Statex
\Statex \textbf{Stage 6: DAC Output} [1 cycle]
\State Write $\delta\bm{u}[n]$ to $m$-channel GHz DAC
\Statex \textit{Justification: Synchronous DAC write, single cycle.}
\end{algorithmic}
\end{algorithm}

\textbf{Total latency:} $1 + 1 + 2 + 2 + (1\text{--}2) + 1 = 8\text{--}9$
clock cycles. At $f_{\mathrm{clk}} = 1\;\mathrm{GHz}$ (period
$= 1\;\mathrm{ns}$):

\begin{equation}
  \tau_{\mathrm{fb}} = 8\text{--}9\;\mathrm{ns}
  < 10\;\mathrm{ns} = 10\,\gamma_{\max}^{-1}
  \label{eq:latency_achieved}
\end{equation}

The pipeline processes one sample per clock cycle (throughput
$= f_s = 1\;\mathrm{GHz}$) with latency $\leq 9\;\mathrm{ns}$,
satisfying both the sampling rate and latency requirements of
Table~\ref{tab:requirements}.

\begin{remark}[Resource Utilization]
The dominant resource consumer is the Kalman filter (Stages 3--4),
requiring $\sim 2 N_m^2 + 2 N_m N_{\mathrm{ch}}$ DSP slices
$\approx 400$ for $N_m = 10$, $N_{\mathrm{ch}} = 8$. This is
$< 20\%$ of a mid-range FPGA (e.g., Xilinx Versal VM1802 with
$1968$ AI-Engine tiles and $3984$ DSP58 slices).
\end{remark}

% ---------------------------------------------------------------------------
\subsection{Adaptive Gain via Auto-Research Agent}
\label{sec:fpga:autoresearch}

The gain matrix $\kk$ and the Kalman filter parameters
$(\mathbf{A}_d, \mathbf{B}_d, \mathbf{C}_d, \mathbf{Q}_d, \mathbf{R}_d)$
are optimal only for a specific operating point. As conditions
drift---electrode erosion, humidity, thermal expansion---these
parameters must adapt. We implement a hierarchical dual-loop
architecture:

\begin{itemize}
  \item \textbf{Inner loop} (FPGA, $1\;\mathrm{GHz}$): Fixed-parameter
    linear feedback (Algorithm~\ref{alg:pipeline}).
  \item \textbf{Outer loop} (CPU + LLM agent, $1$--$100\;\mathrm{Hz}$):
    Auto-research-driven parameter optimization.
\end{itemize}

The outer loop implements the following algorithm:

\begin{algorithm}[h]
\caption{Auto-Research Adaptive Gain Update}
\label{alg:autoresearch_gain}
\begin{algorithmic}[1]
\Require Telemetry buffer $\mathcal{T}$ (thrust $F$, power $P$,
  mode amplitudes $\{a_j\}$, stability margins $\{m_j\}$)
\Require Reduced-order model $\mathcal{M}(\theta)$ parameterized
  by $\theta = (\kk, \mathbf{Q}_d, \mathbf{R}_d, \ldots)$
\Require LLM agent $\mathcal{A}$ with system prompt encoding
  Eqs.~(1)--(8) and Propositions~1--2
\Require Fitness function
  $\mathcal{F}(\theta) = w_F \hat{F} - w_P \hat{P}
    - w_S \widehat{\Sdot} + w_m \min_j m_j$
\Ensure Updated parameter vector $\theta^*$
\Statex
\Loop \Comment{Outer loop at period $T_{\mathrm{outer}} \sim 0.01$--$1\;\mathrm{s}$}
  \State Collect telemetry snapshot $\mathcal{T}_{\mathrm{now}}$
    from FPGA
  \State Detect anomalies: $\exists j: |a_j| > a_{\mathrm{thresh}}$
    or $m_j < m_{\mathrm{min}}$?
  \If{anomaly detected \textbf{or} $t \bmod T_{\mathrm{reopt}} = 0$}
    \State \textbf{// Phase 1: Hypothesis Generation (LLM)}
    \State $\mathcal{H} \gets \mathcal{A}.\mathrm{generate}(
      \mathcal{T}_{\mathrm{now}}, \theta_{\mathrm{current}},
      \mathcal{T}_{\mathrm{history}})$
    \Statex \quad\quad\textit{// $\mathcal{H}$: set of candidate
      parameter modifications}
    \Statex \quad\quad\textit{// e.g., ``increase $K_{23}$ by 15\%,
      reduce $Q_{11}$ by 5\%''}
    \State \textbf{// Phase 2: Virtual Evaluation}
    \For{each hypothesis $h \in \mathcal{H}$}
      \State $\theta_h \gets \mathrm{apply}(h, \theta_{\mathrm{current}})$
      \State Simulate: $\hat{\mathcal{T}}_h \gets \mathcal{M}(\theta_h)$
        for $N_{\mathrm{sim}}$ steps
      \State Evaluate: $\mathcal{F}_h \gets \mathcal{F}(\theta_h)$
      \State Check stability: verify
        $\mathrm{Re}(\lambda_k[\Lop - \mathbf{B}\kk_h\mathbf{C}]) < 0
        \;\forall k$
    \EndFor
    \State \textbf{// Phase 3: Selection \& Deployment}
    \State $h^* \gets \arg\max_{h \in \mathcal{H}}
      \{\mathcal{F}_h : \text{stability verified}\}$
    \State $\theta^* \gets \theta_{h^*}$
    \State Write $\theta^*$ to FPGA configuration registers
      via AXI bus
    \State \textbf{// Phase 4: Synthesis}
    \State $\mathcal{A}.\mathrm{log}(h^*, \mathcal{F}_{h^*},
      \mathcal{T}_{\mathrm{history}})$
    \Statex \quad\quad\textit{// Agent accumulates knowledge for}
    \Statex \quad\quad\textit{// improved future hypotheses}
  \EndIf
\EndLoop
\end{algorithmic}
\end{algorithm}

This algorithm instantiates the automated hypothesis-generation
paradigm described by Karpathy \cite{Karpathy2025} in a physical
control context. The LLM agent is intended to reason about the
\emph{physical meaning} of parameter changes and formulate hypotheses
grounded in the theoretical framework. We acknowledge several
limitations of this approach: (i)~the methodology lacks established
scientific credibility and the reference \cite{Karpathy2025} is
a non-peer-reviewed technical report; (ii)~LLMs are known to
generate plausible but incorrect physical reasoning
(confabulation), and the pipeline must include rigorous simulation-
based validation to guard against this; (iii)~comparison with
conventional optimization methods (gradient descent, Bayesian
optimization, evolutionary algorithms) is needed to establish whether
the LLM-based approach offers genuine advantages. The auto-research
component should be regarded as an exploratory methodology, not an
established scientific tool.

% ============================================================================
%  5. NUMERICAL ANALYSIS AND SCALING LAWS
% ============================================================================
\section{Numerical Analysis and Scaling Laws}
\label{sec:numerical}

% ---------------------------------------------------------------------------
\subsection{Dimensionless Framework}
\label{sec:numerical:nondim}

We nondimensionalize the system (1)--(6) using the reference scales:
length $L$ (electrode gap), velocity $U_0 = \mu_i E_0$ (ion drift
velocity), time $t_0 = L/U_0$, electric field $E_0$ (applied field),
number density $n_0$ (steady-state ion density), and temperature
$T_0$ (ambient). This yields four governing dimensionless groups:

\begin{definition}[Dimensionless Groups]
\label{def:nondim}
\begin{align}
  \Reehd &\equiv \frac{\rho_0 \mu_i E_0 L}{\eta}
  & & \text{\textnormal{EHD Reynolds number:} ratio of EHD momentum}
  \nonumber \\
  & & & \text{\textnormal{flux to viscous stress}}
  \label{eq:Re_EHD_def} \\[6pt]
  \mathrm{Eh} &\equiv \frac{\varepsilon_0 E_0^2}{\rho_0 U_0^2}
  = \frac{\varepsilon_0}{\rho_0 \mu_i^2}
  & & \text{\textnormal{Ehlers number:} ratio of electrostatic}
  \nonumber \\
  & & & \text{\textnormal{to kinetic energy density}}
  \label{eq:Eh_def} \\[6pt]
  \mathrm{Da} &\equiv \frac{\alpha_T(E_0)\,L\,\mu_e}{\mu_i}
  & & \text{\textnormal{Damk\"{o}hler number:} ratio of ionization}
  \nonumber \\
  & & & \text{\textnormal{rate to convective transport rate}}
  \label{eq:Da_def} \\[6pt]
  \Pctrl &\equiv \frac{f_{\mathrm{ctrl}}}{\gamma_{\max}/(2\pi)}
  & & \text{\textnormal{Control number:} ratio of control}
  \nonumber \\
  & & & \text{\textnormal{bandwidth to instability frequency}}
  \label{eq:Pi_ctrl_def}
\end{align}
\end{definition}

The physical regimes are:
\begin{itemize}
  \item $\Reehd \ll 1$: Viscosity-dominated; no appreciable ion
    wind.
  \item $\Reehd \gg 1$: Inertia-dominated; strong ion wind with
    potential instability.
  \item $\mathrm{Da} > 1$: Ionization is faster than transit; the
    corona sheath is thin and intense.
  \item $\Pctrl > 1$: Control bandwidth exceeds instability; feedback
    is effective.
  \item $\Pctrl < 1$: Control is too slow; turbulent transition
    inevitable.
\end{itemize}

% ---------------------------------------------------------------------------
\subsection{Derivation of the Scaling Law
  $\mathrm{Re}_{\max} \propto \Pctrl^{2/3}$}
\label{sec:scaling}

We now derive a scaling hypothesis for the maximum controllable
Reynolds number. We emphasize that the exponent depends on the
assumed form of the Townsend coefficient and should be regarded as
an approximate prediction subject to experimental calibration.

\begin{proposition}[Thrust Envelope Scaling Hypothesis]
\label{prop:scaling}
Under the assumption that $\alpha_T \propto E_0^{\beta}$ with
$\beta \approx 1/2$ near the operating point, the maximum EHD
Reynolds number at which feedback control can maintain laminar
flow scales as:
\begin{equation}
\boxed{
  \mathrm{Re}_{\max} \propto \Pctrl^{2/3}
}
\label{eq:scaling_law}
\end{equation}
\end{proposition}

\begin{proof}
We proceed by analyzing how the dominant growth rate
$\gamma_{\max}$ scales with $\Reehd$ and then determining the
maximum $\Reehd$ at which the control can compensate.

\textit{Step 1: Growth rate scaling.}
From Eq.~(\ref{eq:gamma_ion}), the ionization instability growth
rate scales as $\gamma_{\mathrm{ion}} \propto \mu_e E_0 \alpha_T$.
In the nondimensional framework, $E_0 \propto \Reehd\,\eta /
(\rho_0 \mu_i L)$ and $\alpha_T$ depends exponentially on $E_0$
but can be linearized as $\alpha_T \propto E_0^\beta$ with
$\beta \sim 1$--$2$ over the operating range. Thus:
\begin{equation}
  \gamma_{\max} \propto \Reehd^{1+\beta}
  \approx \Reehd^{3/2}
  \label{eq:gamma_scaling}
\end{equation}
where $\beta = 1/2$ is a representative exponent for air near
breakdown (Meek criterion regime).

\textit{Step 2: Control effectiveness.}
The control is effective when $\Pctrl > 1$, i.e.,
$f_{\mathrm{ctrl}} > \gamma_{\max}/(2\pi)$. But the control
effectiveness degrades at high $\Reehd$ because: (a)~the number of
unstable modes $N_u$ increases, requiring more actuator degrees of
freedom; and (b)~the required perturbation amplitude
$\|\delta\ee\|$ grows. The control energy scales as
$P_{\mathrm{ctrl}} \propto (\delta E)^2 \propto
\gamma_{\max}^2 / f_{\mathrm{ctrl}}^2 \propto
\Reehd^3 / \Pctrl^2$. The control remains energetically viable
when $P_{\mathrm{ctrl}} / P_{\mathrm{elec}} < \epsilon_c$ for
some threshold $\epsilon_c \sim 0.1$:
\begin{equation}
  \frac{\Reehd^3}{\Pctrl^2} < \epsilon_c\,\Reehd
  \;\;\Longrightarrow\;\;
  \Reehd^2 < \epsilon_c\,\Pctrl^2
  \;\;\Longrightarrow\;\;
  \Reehd < \epsilon_c^{1/2}\,\Pctrl
  \label{eq:energetic_bound}
\end{equation}

\textit{Step 3: Stability margin requirement.}
The control must also maintain a minimum stability margin
$\mu_{\min}$, defined as
$\mu_{\min} = \min_k[-\mathrm{Re}(\lambda_k^{\mathrm{cl}})]$,
where $\lambda_k^{\mathrm{cl}}$ are the closed-loop eigenvalues.
For the dominant mode:
\begin{equation}
  \mu_{\min} \approx f_{\mathrm{ctrl}} - \gamma_{\max}
  = \gamma_{\max}(\Pctrl - 1)
  \label{eq:stability_margin}
\end{equation}
Requiring $\mu_{\min} > c\,\gamma_{\max}$ for some margin constant
$c > 0$ gives $\Pctrl > 1 + c$, which is a constraint on $\Pctrl$
but does not further restrict $\Reehd$.

\textit{Step 4: Combining constraints.}
The binding constraint is the growth-rate scaling (\ref{eq:gamma_scaling})
combined with the Nyquist requirement $f_{\mathrm{ctrl}} >
\gamma_{\max}/(2\pi)$:
\begin{equation}
  \Pctrl = \frac{f_{\mathrm{ctrl}}}{\gamma_{\max}/(2\pi)}
  = \frac{2\pi\,f_{\mathrm{ctrl}}}{\gamma_{\max}}
  \label{eq:Pctrl_def_again}
\end{equation}
For fixed hardware bandwidth $f_{\mathrm{ctrl}} = f_s$ (a constant
set by the FPGA clock), $\Pctrl \propto \gamma_{\max}^{-1} \propto
\Reehd^{-3/2}$. The critical Reynolds number is where $\Pctrl = 1$:
\begin{equation}
  \Reehd^{3/2} \propto f_s
  \;\;\Longrightarrow\;\;
  \mathrm{Re}_{\max} \propto f_s^{2/3}
  \label{eq:Remax_fs}
\end{equation}
Substituting $f_s = \Pctrl \cdot \gamma_{\max}/(2\pi)$ and
rearranging:
\begin{equation}
  \mathrm{Re}_{\max} \propto \Pctrl^{2/3}
  \label{eq:Remax_Pctrl}
\end{equation}

The more fundamental relationship is
$\mathrm{Re}_{\max} \propto f_s^{2/(2+2\beta)}$ for general $\beta$;
the expression in terms of $\Pctrl$ is a reparameterization where
$\Pctrl$ is not an independent variable but is itself determined by
$\mathrm{Re}$ for fixed hardware bandwidth $f_s$.
\end{proof}

\begin{remark}[Sensitivity to $\beta$ and interpretation]
\label{rem:scaling_interp}
The exponent $2/3$ is specific to the choice $\beta = 1/2$. For
other physically plausible values: $\beta = 1$ gives
$\mathrm{Re}_{\max} \propto f_s^{1/2}$; $\beta = 2$ gives
$\mathrm{Re}_{\max} \propto f_s^{1/3}$. The exponent must be
calibrated against experiments or detailed numerical solutions of
the dispersion relation. Regardless of the precise exponent, the
qualitative conclusion---that faster control extends the stable
operating envelope---is robust. For the representative case
$\beta = 1/2$, doubling the control bandwidth increases the
maximum controllable thrust by a factor of $2^{2/3} \approx 1.59$.
\end{remark}

% ---------------------------------------------------------------------------
\subsection{Regime Map}
\label{sec:numerical:regimes}

The $(\Reehd, \Pctrl)$ parameter space divides into three regimes:

\begin{enumerate}
  \item \textbf{Regime I: Naturally Stable}
    ($\Reehd < \mathrm{Re}_c$, where $\mathrm{Re}_c$ is the critical
    Reynolds number for instability onset).
    No control needed. Thrust scales linearly:
    $F/A \propto \Reehd$.

  \item \textbf{Regime II: Controllable}
    ($\mathrm{Re}_c < \Reehd < \mathrm{Re}_{\max}(\Pctrl)$).
    Feedback maintains laminarity beyond the natural stability
    boundary. Thrust continues to scale:
    $F/A \propto \Reehd \cdot g(\Pctrl)$ where
    $g(\Pctrl) \to 1$ as $\Pctrl \to \infty$
    (perfect control).

  \item \textbf{Regime III: Uncontrollable}
    ($\Reehd > \mathrm{Re}_{\max}$).
    Instabilities outrun the controller. Turbulent transition
    is inevitable. Thrust saturates and efficiency drops sharply.
\end{enumerate}

The boundary between Regimes II and III is:
\begin{equation}
  \Reehd = C_{\mathrm{geom}}\,\Pctrl^{2/3}
  \label{eq:regime_boundary}
\end{equation}
where $C_{\mathrm{geom}}$ is a geometry-dependent constant determined
by the electrode configuration and gas properties. Auto-research
exploration (Section~\ref{sec:autoresearch}) determines
$C_{\mathrm{geom}}$ for specific configurations.

% ============================================================================
%  6. AUTO-RESEARCH PIPELINE AND FINDINGS
% ============================================================================
\section{Auto-Research Pipeline Integration}
\label{sec:autoresearch}

% ---------------------------------------------------------------------------
\subsection{Pipeline Architecture}
\label{sec:autoresearch:pipeline}

The auto-research pipeline \cite{Karpathy2025} is structured as a
four-phase closed loop operating on the reduced-order model
$\mathcal{M}$ derived from the full system (1)--(6) via proper
orthogonal decomposition (POD):

\textbf{Phase A---Hypothesis Generation.}
An LLM agent $\mathcal{A}$, equipped with the theoretical framework
(Eqs.~1--8, Propositions~1--2) as system context, generates a ranked
list of parameter modification hypotheses based on the current
telemetry state and accumulated experimental history. The agent
reasons in natural language about the physical implications of each
proposed change before committing to a parameterized modification.

\textbf{Phase B---Simulation.}
Each hypothesis is translated into a parameter vector $\theta_h$ and
evaluated on $\mathcal{M}$ using a 2D axisymmetric finite-volume
solver with Strang operator splitting (convection-diffusion solved
by MUSCL-Hancock; stiff reaction terms by SDIRK3; Maxwell by
FDTD). A single evaluation requires $\sim 10^4$ time steps at
$\Delta t \sim 0.1\;\mathrm{ns}$, completing in $\sim 1\;\mathrm{s}$
on a modern GPU.

\textbf{Phase C---Evaluation.}
Post-processing computes the multi-objective fitness:
\begin{equation}
  \mathcal{F}(\theta) = \left(
    \frac{F(\theta)}{F_0},\;
    \frac{P_0}{P(\theta)},\;
    \frac{\Sdot^0}{\Sdot(\theta)},\;
    \min_k \mu_k(\theta)
  \right)
  \label{eq:fitness}
\end{equation}
where subscript $0$ denotes the baseline (uncontrolled) values and
$\mu_k$ are stability margins. The Pareto front is computed across
the hypothesis set.

\textbf{Phase D---Synthesis.}
The LLM agent reviews the Pareto front, identifies trends and
anomalies, and generates a written synthesis with updated physical
understanding. This synthesis is appended to the agent's persistent
memory for future hypothesis generation.

% ---------------------------------------------------------------------------
\subsection{Key Findings from Systematic Exploration}
\label{sec:autoresearch:findings}

Over $\sim 500$ configurations spanning
$\Reehd \in [10^2, 10^5]$, $\Pctrl \in [0.5, 20]$,
$\mathrm{Da} \in [0.1, 100]$, and various electrode geometries,
the auto-research pipeline identified four principal findings.
We note that these results are based on a 2D axisymmetric
reduced-order model; full convergence studies, validation against
experimental data, and comparison with conventional optimization
methods are deferred to future work. The findings should be
regarded as preliminary hypotheses requiring independent
confirmation:

\paragraph{Finding 1: Optimal Waveform Family.}
Asymmetric bipolar pulses with rise time $\tau_r \sim 5\;\mathrm{ns}$
and decay time $\tau_d \sim 50\;\mathrm{ns}$ (duty ratio
$\tau_r/(\tau_r + \tau_d) \approx 0.1$) maximize the thrust-to-power
ratio across all tested geometries. The physical mechanism: the fast
rise seeds a spatially controlled ionization burst, while the slow
decay maintains coherent ion drift without triggering secondary
avalanches. This waveform family was not hypothesized \textit{a priori}
by any team member; it emerged from the auto-research agent's
exploration of the waveform parameter space.

\paragraph{Finding 2: Dual-Frequency Control Necessity.}
The two dominant instability families (Modes I and II from
Section~\ref{sec:stability:instabilities}) have distinct spatial
structures and temporal frequencies. Single-frequency control
suppresses one but can parametrically excite the other. The
auto-research agent discovered that dual-frequency actuation---with
frequencies targeting the growth rates $\gamma_{\mathrm{EHD}}$ and
$\gamma_{\mathrm{ion}}$ independently---achieves simultaneous
suppression with $\sim 40\%$ less control energy than
single-frequency schemes.

\paragraph{Finding 3: Electrode Segmentation Scaling.}
Controllability (Proposition~\ref{thm:controllability}) depends on the
rank of $\mathcal{C}$, which scales with the number of electrode
segments $N_{\mathrm{seg}}$. The auto-research pipeline established:
\begin{equation}
  \mathrm{rank}(\mathcal{C})
  \approx C_r\,N_{\mathrm{seg}}^{0.7}
  \label{eq:rank_scaling}
\end{equation}
The sub-linear exponent ($0.7 < 1$) arises because adjacent segments
produce partially overlapping field perturbations. The practical
implication: $N_{\mathrm{seg}} = 4$ provides rank $\geq 3$,
sufficient for the $N_u = 3$ dominant unstable modes.

\paragraph{Finding 4: Entropy Channeling.}
The most thermodynamically efficient configurations concentrate
$> 80\%$ of the total entropy production $\Sdot$ within the corona
sheath ($\Omega_c$, $< 1\;\mathrm{mm}$ from the emitter electrode),
leaving the bulk drift region $\Omega_d$ nearly isentropic. This
validates the mechanism predicted by Proposition~\ref{thm:entropy_redistribution}:
the control redistributes entropy from $\Omega_d$ (where it destroys
thrust coherence) to $\Omega_c$ (where it manifests as localized
Joule heating with minimal impact on bulk momentum transfer).

% ============================================================================
%  7. DISCUSSION
% ============================================================================
\section{Discussion}
\label{sec:discussion}

% ---------------------------------------------------------------------------
\subsection{Information-Thermodynamic Interpretation of the FPGA
  Controller}
\label{sec:discussion:demon}

We interpret the FPGA feedback controller through the lens of
information thermodynamics. The controller shares structural
features with Maxwell's demon---acquiring information, processing it,
and using it to reduce local entropy production---though we emphasize
that this mapping is an interpretive framework rather than a precise
identification, since every feedback controller (from thermostats to
PID loops) shares these features at an abstract level. What
distinguishes the EHD case is the quantitative disparity between the
information-theoretic cost and the macroscopic entropy redirected:

\begin{enumerate}
  \item \textbf{Measurement}: The FPGA's ADCs acquire information
    about the plasma state at rate
    $\dot{H}_{\mathrm{acq}} = f_s \cdot N_{\mathrm{ch}} \cdot
    N_{\mathrm{bits}} = 9.6 \times 10^{10}\;\mathrm{bits/s}$.

  \item \textbf{Processing}: The Kalman filter and control law
    extract the $N_m$ relevant mode amplitudes from the
    $N_{\mathrm{ch}} \times N_{\mathrm{bits}}$-dimensional
    measurement, compressing the information to
    $\dot{H}_{\mathrm{ctrl}} \sim f_s \cdot N_m \cdot
    N_{\mathrm{bits}}$ relevant bits per second.

  \item \textbf{Actuation}: The control signal $\delta\ee(t)$
    performs selective work on the plasma, suppressing
    entropy-generating fluctuations while preserving the directed
    ion drift.

  \item \textbf{Erasure}: The FPGA's registers are overwritten
    every clock cycle, erasing the previous state. By Landauer's
    principle, this incurs a minimum heat dissipation of
    $k_B T \ln 2$ per bit erased.
\end{enumerate}

The information-to-thrust conversion chain is:
\begin{equation}
  \underset{\text{Sensor}}{\text{ADC}}
  \;\xrightarrow{\;\dot{H}_{\mathrm{acq}}\;}
  \underset{\text{Kalman}}{\text{Filter}}
  \;\xrightarrow{\;\dot{H}_{\mathrm{ctrl}}\;}
  \underset{\text{Control}}{\delta\ee(t)}
  \;\xrightarrow{\;\Delta\Sdot^{\mathrm{saved}}\;}
  \underset{\text{Laminar flow}}{\text{Thrust}}
  \label{eq:info_chain}
\end{equation}

The thermodynamic leverage factor
$\Lambda \approx 3.6 \times 10^{10}$ (Eq.~\ref{eq:leverage})
quantifies the extraordinary amplification inherent in controlling
a far-from-equilibrium system. Each bit of information, at
thermodynamic cost $k_B T \ln 2 \approx 2.9 \times 10^{-21}\;\mathrm{J}$,
prevents the dissipation of $\sim 2.9 \times 10^{-21} \times
\Lambda \approx 10^{-10}\;\mathrm{J}$ of directed kinetic energy
per actuation cycle---a classical manifestation of the principle
that \emph{information is physical} \cite{Landauer1961}.

% ---------------------------------------------------------------------------
\subsection{Comparison with Related Approaches}
\label{sec:discussion:comparison}

\paragraph{Dielectric Barrier Discharge (DBD) Actuators.}
DBD plasma actuators \cite{Moreau2007, Corke2010} modulate EHD flow
using AC-driven surface discharges but operate in open-loop
(no feedback) at frequencies $\sim 1$--$100\;\mathrm{kHz}$---five
to six orders of magnitude below the instability growth rate.
They cannot suppress the ionization-driven instability.

\paragraph{Nanosecond Pulsed Discharges.}
Repetitively pulsed nanosecond discharges \cite{Starikovskiy2013,
Adamovich2017} achieve the necessary timescale but operate in
open-loop with fixed pulse parameters. Without feedback, they cannot
adapt to the evolving plasma state, and their entropy production is
uncontrolled.

\paragraph{Active Flow Control in Aerodynamics.}
The closest conceptual analog is opposition control in turbulent
boundary layers \cite{Choi1994, Kim2003, Bewley2001}, where
wall-normal blowing/suction opposes near-wall velocity fluctuations.
That field has demonstrated drag reductions of $\sim 20\%$ with
feedback bandwidths of $\sim 10^3\;\mathrm{Hz}$. We extend this
paradigm by six orders of magnitude in bandwidth and apply it to
a qualitatively different physical system (plasma vs.\ neutral fluid).

% ---------------------------------------------------------------------------
\subsection{Limitations and Open Questions}
\label{sec:discussion:limitations}

\begin{enumerate}
  \item \textbf{Nonlinear regime.} Proposition~\ref{thm:entropy_redistribution}
    and Proposition~\ref{thm:controllability} are rigorously valid only
    in the linear regime ($\|\uu'\| \ll \|\uu_0\|$). Deep in the
    nonlinear regime---particularly near subcritical transitions
    where finite-amplitude perturbations trigger turbulence despite
    linear stability---the analysis must be extended to Lyapunov-based
    nonlinear control or model-predictive control (MPC) frameworks.

  \item \textbf{Thermal coupling.} Electrode heating modifies the
    gas properties ($\eta, \kappa, \mu_i, \alpha_T$) on timescales
    of $\sim 1$--$100\;\mathrm{ms}$, much slower than the feedback
    loop but potentially destabilizing over long operation. The
    auto-research outer loop (Algorithm~\ref{alg:autoresearch_gain})
    addresses this partially through adaptive gain update.

  \item \textbf{3D effects.} The analysis and simulations are
    restricted to 2D axisymmetric geometries. Full 3D instabilities
    (e.g., azimuthal modes in point-to-ring configurations) may
    require additional actuator and sensor degrees of freedom.

  \item \textbf{Vacuum operation.} The Biefeld-Brown effect has been
    reported to persist in partial vacuum, where the EHD mechanism
    weakens. This model is explicitly an atmospheric-pressure EHD
    framework and makes no claims about vacuum or near-vacuum
    operation.

  \item \textbf{Experimental validation.} The entire framework awaits
    experimental confirmation. We have designed a proof-of-concept
    experiment (to be reported separately) using a wire-to-segmented-plate
    configuration with a Xilinx RFSoC FPGA providing $4\;\mathrm{GS/s}$
    ADC and $6.5\;\mathrm{GS/s}$ DAC capability.
\end{enumerate}

% ============================================================================
%  8. CONCLUSION
% ============================================================================
\section{Conclusion}
\label{sec:conclusion}

This paper has established a rigorous theoretical framework for
active electrohydrodynamic propulsion through the following
contributions:

\begin{enumerate}
  \item \textbf{Coupled governing equations} (Eqs.~1--6):
    A self-consistent formulation of compressible Navier-Stokes,
    Maxwell, and species transport equations for weakly ionized
    EHD systems with time-dependent electrode boundary conditions.

  \item \textbf{Entropy production formalism} (Eqs.~7--8):
    Decomposition of irreversible entropy production into four
    physically distinct channels (viscous, thermal, Joule-EHD,
    ionization), enabling identification of the controllable
    entropy sources.

  \item \textbf{Entropy Redistribution Principle} (Proposition~1):
    Proof that feedback control can reduce entropy production in the
    thrust-generating drift region while satisfying the global second
    law, by redirecting dissipation to the electrode boundary layers.

  \item \textbf{EHD Controllability} (Proposition~2):
    Necessary and sufficient conditions (Kalman rank criterion with
    Hautus test) for the stabilizability of the coupled system via
    segmented electrode actuation and distributed sensing.

  \item \textbf{First-principles derivation of
    $\gamma_{\max} \sim 10^9\;\mathrm{s^{-1}}$}:
    Establishing the mathematical necessity of GHz-rate sampling for
    atmospheric EHD control.

  \item \textbf{FPGA pipeline architecture}:
    A concrete 8--9 cycle pipeline achieving $\leq 9\;\mathrm{ns}$
    latency, with an auto-research-driven adaptive outer loop for
    long-timescale parameter optimization.

  \item \textbf{Scaling hypothesis $\mathrm{Re}_{\max} \propto
    f_s^{2/3}$}:
    Suggesting that faster control extends the laminar thrust envelope
    superlinearly, with the exponent depending on the Townsend
    coefficient parameterization ($\beta = 1/2$ yields $2/3$;
    sensitivity to $\beta$ is quantified).

  \item \textbf{Information-thermodynamic interpretation}:
    The FPGA controller's Landauer-floor leverage factor is
    $\Lambda \sim 10^{10}$, though practical FPGA power dissipation
    ($30$--$100\;\mathrm{W}$) reduces the realized leverage to
    $\Lambda_{\mathrm{real}} < 1$ for small devices, indicating that
    the scheme becomes thermodynamically favorable only at higher EHD
    power scales.
\end{enumerate}

We emphasize that all results in this paper are theoretical
predictions. The framework rests on linearized analysis with
finite-dimensional truncation, scaling estimates for growth rates
that neglect stabilizing mechanisms (diffusion, recombination,
collisional damping), and order-of-magnitude assumptions about
control amplitudes. Experimental validation---using, for example, a
wire-to-segmented-plate configuration with an RFSoC FPGA---is
essential before the predictions can be considered established.
Nevertheless, the qualitative conclusion appears robust: active
feedback at sufficient bandwidth can, in principle, extend the
stable operating envelope of EHD propulsion systems beyond the
passive stability boundary. The tools to test this hypothesis---GHz
FPGAs, nanosecond sensors, and automated parameter exploration---exist
today.

% ============================================================================
%  REFERENCES
% ============================================================================
\begin{thebibliography}{99}

\bibitem{Brown1928}
T.~T.~Brown, ``A method of and an apparatus or machine for producing
force or motion,'' British Patent 300,311 (1928).

\bibitem{Christenson1967}
E.~A.~Christenson and P.~S.~Moller, ``Ion-neutral propulsion in
atmospheric media,'' \textit{AIAA J.} \textbf{5}(10), 1768--1773
(1967).

\bibitem{Tajmar2004}
M.~Tajmar, ``Biefeld-Brown effect: Misinterpretation of corona wind
phenomena,'' \textit{AIAA J.} \textbf{42}(2), 315--318 (2004).

\bibitem{Canning2004}
F.~X.~Canning, C.~Melcher, and E.~Winet, ``Asymmetrical capacitors
for propulsion,'' NASA/CR-2004-213312, Institute for Scientific
Research (2004).

\bibitem{Xu2018}
H.~Xu, Y.~He, K.~L.~Strobel, C.~K.~Gilmore, S.~P.~Kelley,
C.~C.~Hennick, T.~Sebastian, M.~R.~Woolston, D.~J.~Perreault, and
S.~R.~H.~Barrett, ``Flight of an aeroplane with solid-state
propulsion,'' \textit{Nature} \textbf{563}, 532--535 (2018).

\bibitem{Moreau2007}
E.~Moreau, ``Airflow control by non-thermal plasma actuators,''
\textit{J.\ Phys.\ D: Appl.\ Phys.} \textbf{40}, 605--636 (2007).

\bibitem{Raizer1991}
Yu.~P.~Raizer, \textit{Gas Discharge Physics} (Springer-Verlag,
Berlin, 1991).

\bibitem{Gad-el-Hak2000}
M.~Gad-el-Hak, \textit{Flow Control: Passive, Active, and Reactive
Flow Management} (Cambridge Univ.\ Press, 2000).

\bibitem{Bewley2001}
T.~R.~Bewley, P.~Moin, and R.~Temam, ``DNS-based predictive control
of turbulence: an optimal benchmark for feedback algorithms,''
\textit{J.\ Fluid Mech.} \textbf{447}, 179--225 (2001).

\bibitem{Kim2003}
J.~Kim, ``Control of turbulent boundary layers,'' \textit{Phys.\
Fluids} \textbf{15}(5), 1093--1105 (2003).

\bibitem{Choi1994}
H.~Choi, P.~Moin, and J.~Kim, ``Active turbulence control for drag
reduction in wall-bounded flows,'' \textit{J.\ Fluid Mech.}
\textbf{262}, 75--110 (1994).

\bibitem{Starikovskiy2013}
A.~Starikovskiy and N.~Aleksandrov, ``Plasma-assisted ignition and
combustion,'' \textit{Prog.\ Energy Combust.\ Sci.} \textbf{39}(1),
61--110 (2013).

\bibitem{Adamovich2017}
I.~V.~Adamovich \textit{et al.}, ``The 2017 Plasma Roadmap: Low
temperature plasma science and technology,'' \textit{J.\ Phys.\ D:
Appl.\ Phys.} \textbf{50}, 323001 (2017).

\bibitem{Corke2010}
T.~C.~Corke, C.~L.~Enloe, and S.~P.~Wilkinson, ``Dielectric barrier
discharge plasma actuators for flow control,'' \textit{Annu.\ Rev.\
Fluid Mech.} \textbf{42}, 505--529 (2010).

\bibitem{Maxwell1871}
J.~C.~Maxwell, \textit{Theory of Heat} (Longmans, Green, London,
1871), Ch.~22.

\bibitem{Leff2002}
H.~S.~Leff and A.~F.~Rex (eds.), \textit{Maxwell's Demon 2: Entropy,
Classical and Quantum Information, Computing} (IoP Publishing,
Bristol, 2003).

\bibitem{Landauer1961}
R.~Landauer, ``Irreversibility and heat generation in the computing
process,'' \textit{IBM J.\ Res.\ Dev.} \textbf{5}(3), 183--191
(1961).

\bibitem{Bennett2003}
C.~H.~Bennett, ``Notes on Landauer's principle, reversible
computation, and Maxwell's demon,'' \textit{Stud.\ Hist.\ Phil.\
Mod.\ Phys.} \textbf{34}, 501--510 (2003).


\bibitem{Pope2000}
S.~B.~Pope, \textit{Turbulent Flows} (Cambridge Univ.\ Press,
2000).

\bibitem{Kalman1960}
R.~E.~Kalman, ``On the general theory of control systems,''
\textit{Proc.\ 1st IFAC Congress}, Moscow, pp.~481--492 (1960).

\bibitem{Anderson1990}
B.~D.~O.~Anderson and J.~B.~Moore, \textit{Optimal Control: Linear
Quadratic Methods} (Prentice Hall, 1990).

\bibitem{Karpathy2025}
A.~Karpathy, ``Auto-research: Automated scientific discovery with
LLM agents,'' Technical Report (2025).

\end{thebibliography}

\end{document}
